% concatenated from main_23_.tex
\documentclass[11pt]{amsart}

\usepackage{amsmath,amsfonts,amssymb,amsthm,mathrsfs}
\usepackage{hyperref}
\usepackage{enumitem}
%\usepackage{refcheck}
\hypersetup{
    colorlinks=true,
    linkcolor=blue,
    citecolor=blue,
    urlcolor=blue
}
\newtheorem{definition}{Definition}[section]
\newtheorem{theorem}[definition]{Theorem}
\newtheorem{proposition}[definition]{Proposition}
\newtheorem{remark}[definition]{Remark}
\title[Spherical Milnor Spaces and Topological Structures]{
Spherical Milnor Spaces II: Projective Quotients and Higher Topological Structures
}

\author{Jean-Pierre Magnot}

\address{{SFR MATHSTIC, LAREMA, Universit\'e d’Angers, 2 Bd Lavoisier, 
49045 Angers cedex 1, France;  Lyc\'ee Jeanne d'Arc, 40 avenue de Grande Bretagne, 63000 Clermont-Ferrand, 
France}; Lepage Research Institute, 17 novembra 1, 081 16 Presov, Slovakia}
\email{\small magnot@math.cnrs.fr; jean-pierr.magnot@ac-clermont.fr; jp.magnot@gmail.com}


\subjclass[2020]{55R10, 57R20, 58B25, 53C27}

\keywords{
Diffeology, Milnor classifying space, projective structures, 
principal bundles, $\mathbb{Z}_2$-actions, non-abelian gerbes, higher bundles
}

\begin{document}

\begin{abstract}
We introduce a spherical variant of Milnor's classifying construction for diffeological groups, based on quadratic normalization of barycentric coordinates. This construction gives rise to a contractible diffeological space endowed with commuting actions of a group $G$ and of $\mathbb{Z}_2$, leading naturally to a hierarchy of quotient spaces.

We investigate the topological and geometric properties of these quotients, including a projective model and a double quotient space which encodes twisted and higher structures. In particular, we show that this framework provides a natural setting for the study of principal bundles with $\mathbb{Z}_2$-twists, and leads to obstruction classes in low-degree cohomology.

The construction is further related to non-abelian gerbes and higher bundles, providing a bridge between diffeological geometry, classifying space theory, and higher topological structures.
\end{abstract}

\maketitle
\tableofcontents
\section{Introduction}

The construction of classifying spaces is one of the central organizing
principles of modern topology and geometry. Since Milnor's construction of
universal bundles \cite{Milnor1956I,Milnor1956II}, classifying spaces have
provided a natural framework for the study of principal bundles,
characteristic classes, and homotopical invariants. The later developments of
Segal, Dold, Rector, Dupont, Kamber--Tondeur and Karoubi placed this circle of
ideas in a broader homotopical, cohomological and \(K\)-theoretic framework
\cite{Segal1968,Dold1963,Rector1971,Dupont1978,KamberTondeur1975,Karoubi1978}.
In particular, the interaction between classifying spaces, local coefficient
systems, twisted characteristic classes and generalized cohomology theories is
now a standard but still very active part of global geometry.

The present paper is the second part of a program devoted to spherical
versions of Milnor classifying spaces. In Part~I \cite{Magnot-MilnorI}, we
introduced the spherical Milnor model
\[
E^{\mathrm{sph}}(G),
\]
defined by replacing the usual barycentric relation
\[
\sum_i t_i=1,\qquad t_i\geq0,
\]
by the quadratic normalization
\[
\sum_i x_i^2=1.
\]
This change has substantial geometric consequences. The resulting space is not
merely a reparametrization of Milnor's original model: its diffeology differs
from the Milnor--Watts diffeology because no stationarity condition is imposed
at the hypersurfaces \(x_i=0\). The diffeology of Milnor's classifying space
and the role of boundary-adapted plots were studied in \cite{MagnotWatts2017},
within the broader framework of diffeology developed by Souriau and
Iglesias-Zemmour \cite{Souriau1980,IglesiasZemmour2013}. The present paper
builds on that analysis and focuses on the topological, gerbe-theoretic, and
\(K\)-theoretic consequences of the spherical construction.

Diffeology provides a convenient setting for generalized smooth spaces beyond
manifolds. It is flexible enough to include quotients, mapping spaces,
singular spaces, and infinite-dimensional examples, while retaining a usable
notion of smooth parametrization. Foundational and comparative aspects of
diffeology appear in \cite{Hector1995,Stacey2011,Watts2012Thesis,GMW2023}.
The \(D\)-topology and homotopy theory of diffeological spaces were developed
in \cite{ChristensenSinnamonWu2014,ChristensenWu2014Homotopy}, and smooth
classifying spaces were studied in \cite{ChristensenWu2021Classifying}. The
tangent, exterior, and vector-space aspects of diffeology play a role in the
background of this work and are treated in
\cite{ChristensenWu2016Tangent,ChristensenWu2017Filtered,ChristensenWu2019Vector,ChristensenWu2022Exterior}.
Vector pseudo-bundles, pushforwards, and diffeological connections provide the
appropriate framework for the geometric constructions used in Part~I
\cite{Pervova2016,Pervova2017,Pervova2018,Wu2023,MagnotConnections}. Related
views on smooth generalized spaces and infinite-dimensional calculus are also
found in convenient analysis and synthetic differential geometry
\cite{KrieglMichor1997,Kock2006}. The example of the irrational torus shows
how diffeology may interact with noncommutative geometric phenomena
\cite{IglesiasLaffineurZemmour2020}.

The main aim of the present article is to study the topological structures
naturally attached to the spherical Milnor construction. A key feature is the
presence of a canonical \(\mathbb Z_2\)-symmetry
\[
(x_i,g_i)\mapsto (-x_i,g_i),
\]
which leads to projective and twisted quotients. When this symmetry interacts
with an action
\[
\sigma:\mathbb Z_2\to \operatorname{Aut}(G),
\]
one obtains semidirect products
\[
H_\sigma=G\rtimes_\sigma\mathbb Z_2
\]
and corresponding twisted spherical classifying spaces
\[
B^{\mathrm{sph}}_\sigma(G).
\]
The associated transition functions naturally split into a
\(\mathbb Z_2\)-valued cocycle and a \(G\)-valued cocycle twisted by the local
system determined by the first one. This leads to a non-abelian cocycle
picture, closely related to the classical theory of non-abelian cohomology
\cite{Giraud1971} and to higher-categorical approaches to non-abelian
topology \cite{BrownHiggins1981,BrownHigginsSivera2011,BaezLauda2004}.

A central part of the paper is devoted to lifting problems and gerbes. Given
an extension
\[
1\to A\to \widehat H_{\widehat\sigma}\to H_\sigma\to 1,
\]
a principal \(H_\sigma\)-bundle has, in general, an obstruction to being lifted
to a principal \(\widehat H_{\widehat\sigma}\)-bundle. In the presence of the
\(\mathbb Z_2\)-twist, this obstruction lives not in an ordinary coefficient
system but in a local system
\[
A_\alpha
\]
determined by the projective class
\[
\alpha\in H^1(X;\mathbb Z_2).
\]
The corresponding lifting gerbe is banded by \(A_\alpha\), and its class lies
in
\[
H^2(X;A_\alpha).
\]
This point of view connects the present construction with the classical and
modern theories of gerbes, bundle gerbes, and higher stacks
\cite{Breen1994,Breen2008,BreenMessing2005,Murray1996,Stevenson2000,Waldorf2007,NikolausWaldorf2013}.
The language of stacks and orbifold-like quotient structures is also relevant
to the background of these constructions \cite{MoerdijkPronk1997,Noohi2005,Noohi2012}.
The higher bundle and \(\infty\)-geometric viewpoint developed in
\cite{NikolausSchreiberStevenson2012,Schreiber2013,FiorenzaSchreiberStasheff2007}
provides a natural conceptual extension of the present picture. In particular,
the relation between differential cocycles, higher principal bundles and
higher characteristic classes motivates the interpretation of the obstruction
classes introduced below. Waldorf's transgression-regression methods and
models for string \(2\)-groups also provide a useful comparison point
\cite{Waldorf2010}, while Wockel's work relates such questions to classifying
spaces of principal \(2\)-bundles \cite{Wockel2008}.

When \(A=U(1)\), the lifting gerbe determines a twisted
Dixmier--Douady class. Because of the \(\mathbb Z_2\)-local system, this class
naturally belongs to
\[
H^3(X;\mathbb Z_\alpha),
\]
rather than to ordinary integral cohomology. This gives a twisted analogue of
the usual Dixmier--Douady invariant of a \(U(1)\)-banded gerbe. Such classes
are closely related to the bundle-theoretic interpretation of continuous trace
algebras and twisted \(K\)-theory \cite{Rosenberg1989,AtiyahSegal2004}.
The graded Brauer-group approach of Donovan and Karoubi is particularly
relevant here, since the spherical construction naturally combines a degree-one
\(\mathbb Z_2\)-twist with a gerbe-type degree-three twist
\cite{DonovanKaroubi1970}. The present paper therefore associates to the
spherical construction a twisting datum
\[
(\alpha,\operatorname{DD}_\alpha)
\]
and studies the corresponding twisted \(K\)-groups. These recover the usual
gerbe-twisted \(K\)-theory when \(\alpha=0\), and they encode local
coefficient phenomena when \(\alpha\neq0\).

Another motivation comes from operator-theoretic defects. In Part~I, Hodge-
and Dirac-type structures were constructed on the spherical Milnor model under
appropriate finite-rank or analytic assumptions. Such structures need not
descend to the quotient unless they are compatible with the group action. The
failure of descent defines a defect. In favorable restricted Hilbert settings,
this defect can be related to Schwinger cocycles, central extensions and
lifting gerbes. The classical model for this phenomenon is the theory of loop
groups and restricted Grassmannians \cite{PressleySegal1986}, together with
the appearance of gerbes and anomalies in index theory and Hamiltonian
quantization \cite{CareyMickelssonMurray1997,Freed1995}. Abelian and
non-abelian extensions of infinite-dimensional Lie groups provide the natural
background for these constructions \cite{Neeb2004,Neeb2007}. In the present
article, we formulate the Schwinger-cocycle interpretation carefully as a
conditional construction: it requires a polarization, restricted operators, a
well-defined Schwinger cocycle, and an integrability condition to pass from
Lie algebra cocycles to group extensions.

The paper is organized as follows. Section~\ref{sec:prelim} recalls the
diffeological Milnor and spherical Milnor constructions needed in the sequel.
We emphasize the difference between the simplicial and spherical diffeologies,
the role of the projective \(\mathbb Z_2\)-quotient, and the relation with
Part~I. The next sections study strata, smoothness criteria, and the behavior
of forms and maps with respect to the final diffeology. We then introduce
twisted spherical classifying spaces associated with semidirect products
\[
G\rtimes_\sigma \mathbb Z_2
\]
and describe their transition cocycles. The subsequent sections develop the
lifting gerbe associated with an extension of \(H_\sigma\), prove that its
class lies in the twisted group
\[
H^2(X;A_\alpha),
\]
and identify the corresponding twisted Dixmier--Douady invariant when
\(A=U(1)\). We then discuss the associated twists of \(K\)-theory. Finally, we
analyze the descent of geometric structures, define Hodge and Dirac defects
under appropriate analytic assumptions, and relate restricted equivariance
defects to Schwinger cocycles, central extensions and gerbes.

The main contribution of this article is therefore to show that the spherical
Milnor construction does not merely provide an alternative model for
classifying spaces. It naturally produces:
\[
\alpha\in H^1(X;\mathbb Z_2),
\qquad
\beta\in H^2(X;A_\alpha),
\qquad
\operatorname{DD}_\alpha\in H^3(X;\mathbb Z_\alpha),
\]
together with associated twisted \(K\)-theory classes and defect gerbes. In
this sense, the spherical model provides a unified diffeological framework
linking projective quotients, non-abelian cocycles, gerbes, twisted
Dixmier--Douady classes, and operator-theoretic anomalies.
\section{Preliminaries on diffeological Milnor constructions}
\label{sec:prelim}

We briefly recall the constructions and notions developed in Part~I \cite{Magnot-MilnorI} that will
be used throughout the present article. General references on diffeology are
\cite{IglesiasZemmour2013,GMW2023}.

\subsection{Diffeological spaces and quotient constructions}

A diffeological space is a set endowed with a family of parametrizations
(called plots) stable under smooth reparametrization and satisfying locality
and covering conditions. A map between diffeological spaces is smooth if it
pulls back plots to plots.

Given a family of maps
\[
f_i:X_i\to X,
\]
the final diffeology on \(X\) is the finest diffeology making all \(f_i\)
smooth. Quotient diffeologies are special cases of final diffeologies.

A diffeological group is a group equipped with a compatible diffeology for
which multiplication and inversion are smooth.

\subsection{Milnor classifying spaces}

Let \(G\) be a diffeological group.

The Milnor total space is obtained from
\[
\widetilde E(G)
=
\left\{
(t_i,g_i)
\;\middle|\;
t_i\ge0,\ 
\sum_i t_i=1,\ 
g_i\in G,
\text{ finite support}
\right\}
\]
by quotienting with respect to the relation
\[
(t_i,g_i)\sim(t_i,g_i')
\quad\Longleftrightarrow\quad
g_i=g_i'
\text{ whenever }t_i\neq0.
\]

The resulting quotient space is denoted
\[
E(G),
\]
and the classifying space is
\[
B(G)=E(G)/G,
\]
where \(G\) acts diagonally on the group coordinates.

The diffeology on \(E(G)\) is the quotient of the final diffeology generated
by the local models
\[
\Delta_I\times G^I,
\]
where
\[
\Delta_I
=
\left\{
(t_i)_{i\in I}
\;\middle|\;
t_i\ge0,\ 
\sum_{i\in I}t_i=1
\right\}.
\]

In the Milnor--Watts framework \cite{MagnotWatts2017}, the simplices
\(\Delta_I\) are equipped with a boundary-adapted diffeology forcing plots to
be stationary when some coordinate vanishes.

The bundle
\[
E(G)\to B(G)
\]
is a universal diffeological principal \(G\)-bundle for \(D\)-numerable
bundles.

\subsection{The spherical Milnor model}

The spherical Milnor construction replaces the affine condition
\[
\sum_i t_i=1
\]
by the quadratic normalization
\[
\sum_i x_i^2=1.
\]

More precisely,
\[
\widetilde E^{\mathrm{sph}}(G)
=
\left\{
(x_i,g_i)
\;\middle|\;
\sum_i x_i^2=1,\ 
g_i\in G,
\text{ finite support}
\right\},
\]
and the quotient relation is
\[
(x_i,g_i)\sim(x_i,g_i')
\quad\Longleftrightarrow\quad
g_i=g_i'
\text{ whenever }x_i\neq0.
\]

The spherical Milnor space is
\[
E^{\mathrm{sph}}(G)
=
\widetilde E^{\mathrm{sph}}(G)/\sim.
\]

Its diffeology is the quotient of the final diffeology generated by the maps
\[
S_I\times G^I
\longrightarrow
E^{\mathrm{sph}}(G),
\]
where
\[
S_I
=
\left\{
(x_i)_{i\in I}
\;\middle|\;
\sum_{i\in I}x_i^2=1
\right\}
\]
is equipped with the subset diffeology inherited from \(\mathbb R^I\).

Contrary to the simplicial Milnor model, no stationarity condition is imposed
at the loci \(x_i=0\).

\subsection{Comparison of the two models}

There exists a smooth embedding
\[
\iota:E(G)\longrightarrow E^{\mathrm{sph}}(G),
\qquad
(t_i,g_i)\mapsto(\sqrt{t_i},g_i),
\]
because the Milnor diffeology forces stationarity at \(t_i=0\).

Conversely, the map
\[
(x_i,g_i)\mapsto(x_i^2,g_i)
\]
is generally not smooth from the spherical model to the simplicial one.

Therefore the two diffeological constructions are not equivalent. This
difference is responsible for the distinct geometric properties developed in
Part~I.

\subsection{Projective quotient}

The spherical model carries a natural \(\mathbb Z_2\)-action:
\[
\varepsilon\cdot(x_i,g_i)
=
(\varepsilon x_i,g_i),
\qquad
\varepsilon=\pm1.
\]

This action commutes with the \(G\)-action and defines the projective quotient
\[
E^{\mathrm{proj}}(G)
=
E^{\mathrm{sph}}(G)/\mathbb Z_2.
\]

Similarly, one defines the spherical classifying space
\[
B^{\mathrm{sph}}(G)
=
E^{\mathrm{sph}}(G)/(G\times\mathbb Z_2).
\]

\subsection{Differential forms and horizontal geometry}

Differential forms on a diffeological space are defined plotwise: a
\(k\)-form associates to each plot
\[
p:U\to X
\]
a smooth differential form on \(U\), compatible with smooth
reparametrizations.

In Part~I, the spherical Milnor model was endowed with:
\begin{itemize}
\item an effective horizontal pseudo-bundle
\[
\mathcal H
=
TE^{\mathrm{sph}}(G)/\mathcal K;
\]
\item a non-degenerate horizontal metric;
\item horizontal differential forms;
\item Hodge-type operators;
\item Clifford structures and Dirac-type operators.
\end{itemize}

The present article focuses on the topological and homotopical aspects of
these constructions.

\section{Diffeological smoothness and gluing} \label{sec:glue}
In this section we describe more precisely the basic definitions and the properties of the diffeology setteled in \cite{Magnot-MilnorI} on the spherical classifying space. 
\subsection{Final diffeology and plots}

Recall that $E^{\mathrm{sph}}(G)$ is endowed with the quotient of the final diffeology generated by the maps
\[
\psi_I : S_I \times G^I \longrightarrow \widetilde E^{\mathrm{sph}}(G),
\]
followed by the projection to $E^{\mathrm{sph}}(G)$.

\begin{proposition}[Characterization of plots]
A map $p : U \to E^{\mathrm{sph}}(G)$ is a plot if and only if for every $u \in U$, there exist:
\begin{itemize}
\item an open neighborhood $V \subset U$ of $u$,
\item a finite subset $I \subset \mathbb N$,
\item a smooth map $\widetilde p : V \to S_I \times G^I$,
\end{itemize}
such that
\[
p|_V = \psi_I \circ \widetilde p.
\]
\end{proposition}

\begin{proof}
This is the standard characterization of the final diffeology, combined with the quotient diffeology.
\end{proof}

\subsection{Local finiteness}

\begin{proposition}
Every plot $p : U \to E^{\mathrm{sph}}(G)$ is locally supported on a fixed finite index set.
\end{proposition}

\begin{proof}
By definition of the final diffeology, every point admits a neighborhood factoring through some $\psi_I$ with $I$ finite.
\end{proof}

\begin{remark}
Thus all local computations reduce to finite-dimensional models.
\end{remark}

\subsection{Gluing of plots}

\begin{proposition}[Sheaf condition]
Let $(U_\alpha)_\alpha$ be an open cover of $U$, and let $p_\alpha : U_\alpha \to E^{\mathrm{sph}}(G)$ be plots such that
\[
p_\alpha = p_\beta \quad \text{on } U_\alpha \cap U_\beta.
\]
Then the map $p : U \to E^{\mathrm{sph}}(G)$ defined by $p|_{U_\alpha} = p_\alpha$ is a plot.
\end{proposition}

\begin{proof}
This is one of the defining axioms of a diffeology.
\end{proof}

\subsection{Strata and local models}

\begin{definition}
For each finite \(I\subset \mathbb N\), define
\[
E_I:=\psi_I(S_I\times G^I)\subset E^{\mathrm{sph}}(G).
\]
We endow \(E_I\) with the subset diffeology induced from
\(E^{\mathrm{sph}}(G)\).
\end{definition}

\begin{proposition}
For every finite \(I\subset\mathbb N\), the map
\[
\psi_I:S_I\times G^I\longrightarrow E_I
\]
is a subduction.
\end{proposition}

\begin{proof}
By definition, \(E_I\) is the image of \(\psi_I\). Hence \(\psi_I\) is
surjective.

The diffeology on \(E^{\mathrm{sph}}(G)\) is the final diffeology generated by
the maps
\[
\psi_J:S_J\times G^J\to E^{\mathrm{sph}}(G).
\]
The subset diffeology on \(E_I\) consists of those plots of
\(E^{\mathrm{sph}}(G)\) whose image lies in \(E_I\).

Since \(\psi_I\) is one of the generating maps, every plot of
\(S_I\times G^I\) is sent by \(\psi_I\) to a plot of
\(E^{\mathrm{sph}}(G)\), hence to a plot of \(E_I\).

Conversely, let \(p:U\to E_I\) be a plot for the subset diffeology. Then
\(p\) is locally a plot of \(E^{\mathrm{sph}}(G)\), hence locally factors
through some
\[
\psi_J:S_J\times G^J\to E^{\mathrm{sph}}(G).
\]
Since the image of \(p\) is contained in \(E_I\), this local factorization may
be regarded, after replacing representatives by equivalent ones, as a local
factorization through the quotient image \(E_I\). Thus the subset diffeology
on \(E_I\) coincides with the push-forward diffeology induced by
\(\psi_I\). Therefore \(\psi_I\) is a subduction.
\end{proof}

\begin{remark}
The map \(\psi_I\) is generally not injective. Indeed, if \(x_i=0\), then the
corresponding group coordinate \(g_i\) is irrelevant in the quotient relation.
\end{remark}

\subsection{Smoothness criterion}

\begin{proposition}
Let \(Y\) be a diffeological space. A map
\[
f:E^{\mathrm{sph}}(G)\to Y
\]
is smooth if and only if, for every finite \(I\subset\mathbb N\), the
composition
\[
f\circ \psi_I:S_I\times G^I\to Y
\]
is smooth.
\end{proposition}

\begin{proof}
Assume first that \(f\) is smooth. Since each \(\psi_I\) is smooth by
construction of the final diffeology, the composition
\[
f\circ\psi_I
\]
is smooth.

Conversely, assume that every \(f\circ\psi_I\) is smooth. Let
\[
p:U\to E^{\mathrm{sph}}(G)
\]
be a plot. Since the diffeology on \(E^{\mathrm{sph}}(G)\) is final with
respect to the family \((\psi_I)_I\), for every \(u\in U\) there exists an open
neighborhood \(V\subset U\), a finite set \(I\), and a plot
\[
\widetilde p:V\to S_I\times G^I
\]
such that
\[
p|_V=\psi_I\circ\widetilde p.
\]
Then
\[
f\circ p|_V
=
f\circ\psi_I\circ\widetilde p,
\]
which is smooth by assumption. Since smoothness of plots is local on the
domain, \(f\circ p\) is a plot of \(Y\). Hence \(f\) is smooth.
\end{proof}

\subsection{Behavior across strata}

The following point is important. The spherical model does not impose
stationarity at the hypersurfaces \(x_j=0\). Therefore a plot may cross a
boundary stratum transversally.

\begin{remark}
It is not true in general that if a coordinate \(x_j\) vanishes at a point of a
plot, then the plot locally factors through a smaller stratum. For instance,
on a sphere \(S_I\), a smooth curve may satisfy \(x_j(u_0)=0\) while
\(x_j(u)\neq0\) for \(u\neq u_0\). Thus the spherical diffeology allows
genuine transverse crossing of strata.
\end{remark}

The correct statement is the following.

\begin{proposition}
Let
\[
p:U\to E^{\mathrm{sph}}(G)
\]
be a plot. Suppose that locally
\[
p=\psi_I\circ \widetilde p
\]
with
\[
\widetilde p=(x_i,g_i)_{i\in I}:U\to S_I\times G^I.
\]
If, on an open subset \(V\subset U\), one has
\[
x_j|_V\equiv0
\]
for some \(j\in I\), then \(p|_V\) factors through the smaller model
\[
S_{I\setminus\{j\}}\times G^{I\setminus\{j\}}.
\]
\end{proposition}

\begin{proof}
Assume that \(x_j(v)=0\) for every \(v\in V\). Since
\[
\sum_{i\in I}x_i(v)^2=1,
\]
we have
\[
\sum_{i\in I\setminus\{j\}}x_i(v)^2=1
\]
for every \(v\in V\). Hence the remaining coordinates define a smooth map
\[
(x_i)_{i\in I\setminus\{j\}}:V\to S_{I\setminus\{j\}}.
\]

Similarly, by deleting the inactive group coordinate \(g_j\), we obtain a
smooth map
\[
\widetilde p_j:V\to S_{I\setminus\{j\}}\times G^{I\setminus\{j\}}.
\]

Now compare the images of \(\widetilde p|_V\) and \(\widetilde p_j\) in
\(E^{\mathrm{sph}}(G)\). Since the coordinate \(x_j\) is identically zero on
\(V\), the equivalence relation defining \(E^{\mathrm{sph}}(G)\) makes the
corresponding group coordinate \(g_j\) irrelevant. Therefore
\[
\psi_I\circ\widetilde p|_V
=
\psi_{I\setminus\{j\}}\circ\widetilde p_j.
\]
Thus \(p|_V\) factors through the smaller model.
\end{proof}

\begin{proposition}
A plot factors locally through a lower stratum precisely on those open subsets
where the corresponding spherical coordinates vanish identically.
\end{proposition}

\begin{proof}
The preceding proposition proves the sufficient condition. Conversely, if a
plot factors through \(S_{I\setminus\{j\}}\times G^{I\setminus\{j\}}\), then
the coordinate \(x_j\) is identically zero in that local representation.
\end{proof}

\begin{remark}
Thus the strata are not separated by a stationarity condition. The spherical
model allows plots to cross from one support type to another. This is a major
difference from the Milnor--Watts simplicial diffeology, where boundary
stationarity plays a crucial role.
\end{remark}
\subsection{Smooth functions}

\begin{definition}
A function
\[
f:E^{\mathrm{sph}}(G)\to\mathbb R
\]
is called smooth if for every plot
\[
p:U\to E^{\mathrm{sph}}(G),
\]
the composition
\[
f\circ p:U\to\mathbb R
\]
is smooth in the ordinary sense.
\end{definition}

\begin{proposition}
A function
\[
f:E^{\mathrm{sph}}(G)\to\mathbb R
\]
is smooth if and only if, for every finite subset
\[
I\subset\mathbb N,
\]
the composition
\[
f\circ\psi_I:S_I\times G^I\to\mathbb R
\]
is smooth.
\end{proposition}

\begin{proof}
Assume first that \(f\) is smooth. By construction of the diffeology on
\(E^{\mathrm{sph}}(G)\), each map
\[
\psi_I:S_I\times G^I\to E^{\mathrm{sph}}(G)
\]
is smooth. Therefore the composition
\[
f\circ\psi_I
\]
is smooth for every finite \(I\).

Conversely, assume that each
\[
f\circ\psi_I
\]
is smooth. Let
\[
p:U\to E^{\mathrm{sph}}(G)
\]
be a plot. Since the diffeology on \(E^{\mathrm{sph}}(G)\) is the final
diffeology generated by the family \((\psi_I)_I\), for every point
\(u\in U\) there exist:
\begin{itemize}
\item an open neighborhood \(V\subset U\) of \(u\),
\item a finite subset \(I\subset\mathbb N\),
\item and a smooth map
\[
\widetilde p:V\to S_I\times G^I
\]
such that
\[
p|_V=\psi_I\circ\widetilde p.
\]
\end{itemize}

Hence
\[
f\circ p|_V
=
f\circ\psi_I\circ\widetilde p.
\]

By assumption,
\[
f\circ\psi_I
\]
is smooth, and since \(\widetilde p\) is smooth, the composition
\[
f\circ p|_V
\]
is smooth.

Since smoothness is local on the domain, it follows that
\[
f\circ p
\]
is smooth on all of \(U\). Therefore \(f\) is smooth as a diffeological map.
\end{proof}

\subsection{Summary}

The spherical Milnor model
\[
E^{\mathrm{sph}}(G)
\]
is equipped with a diffeological structure characterized by the following
properties:
\begin{itemize}
\item every plot locally factors through a finite-dimensional model
\[
S_I\times G^I;
\]
\item the diffeology is generated by the family of maps
\[
\psi_I:S_I\times G^I\to E^{\mathrm{sph}}(G);
\]
\item smooth maps and smooth functions are completely characterized by their
pull-backs through the generators \(\psi_I\);
\item the space possesses a stratified structure governed by the vanishing of
the spherical coordinates;
\item contrary to the simplicial Milnor model, plots may cross strata
transversally, since no stationarity condition is imposed at the loci
\[
x_i=0.
\]
\end{itemize}

These properties will be used repeatedly in the sequel when constructing
geometric, topological, and analytical structures on the spherical Milnor
space.
\section{Projective quotients and \texorpdfstring{$\mathbb Z_2$}{Z2}-twisted classifying structures}
\label{sec:projective-topology}

\subsection{The antipodal action}

Recall that the spherical Milnor space is
\[
E^{\mathrm{sph}}(G)
=
\widetilde E^{\mathrm{sph}}(G)/\sim,
\]
where
\[
(x_i,g_i)\sim (x_i,g_i')
\quad\Longleftrightarrow\quad
g_i=g_i' \ \text{whenever } x_i\neq0.
\]

\begin{definition}
The antipodal action of \(\mathbb Z_2=\{\pm1\}\) on \(E^{\mathrm{sph}}(G)\) is defined by
\[
\varepsilon\cdot [(x_i,g_i)]
=
[(\varepsilon x_i,g_i)].
\]
\end{definition}

\begin{proposition}
The antipodal action of \(\mathbb Z_2\) on \(E^{\mathrm{sph}}(G)\) is well-defined, smooth and free.
\end{proposition}

\begin{proof}
We first prove that the action is well-defined. Suppose that
\[
(x_i,g_i)\sim (x_i,g_i').
\]
Then \(g_i=g_i'\) for every \(i\) such that \(x_i\neq0\). Since
\[
\varepsilon x_i\neq0
\quad\Longleftrightarrow\quad
x_i\neq0,
\]
we also have \(g_i=g_i'\) for every \(i\) such that \(\varepsilon x_i\neq0\). Hence
\[
(\varepsilon x_i,g_i)\sim(\varepsilon x_i,g_i'),
\]
so the action is well-defined on the quotient.

Smoothness follows from the definition of the final diffeology. On each finite model
\[
S_I\times G^I,
\]
the action is given by
\[
(x_i,g_i)_{i\in I}
\longmapsto
(\varepsilon x_i,g_i)_{i\in I},
\]
which is smooth since \(S_I\subset \mathbb R^I\) carries the subset diffeology.

Finally, suppose that \(-1\) fixes a point \([(x_i,g_i)]\). Then
\[
[(x_i,g_i)]=[(-x_i,g_i)].
\]
The equivalence relation does not change the spherical coordinates, hence this implies
\[
x_i=-x_i
\]
for every \(i\). Therefore \(x_i=0\) for all \(i\), contradicting
\[
\sum_i x_i^2=1.
\]
Thus the action is free.
\end{proof}

\begin{definition}
The projective spherical Milnor space is
\[
E^{\mathrm{proj}}(G)
:=
E^{\mathrm{sph}}(G)/\mathbb Z_2,
\]
equipped with the quotient diffeology.
\end{definition}

\subsection{Contractibility of the spherical total space}

We now prove that \(E^{\mathrm{sph}}(G)\) is contractible. This result is essential: it is the topological reason why projective quotients of \(E^{\mathrm{sph}}(G)\) carry universal \(\mathbb Z_2\)-type classes.

\begin{theorem}
The diffeological space \(E^{\mathrm{sph}}(G)\) is smoothly contractible.
\end{theorem}

\begin{proof}
Let
\[
e_1=(1,0,0,\ldots)
\]
be the first coordinate vector, and choose the constant group label \(1_G\) on all inactive coordinates. We construct a smooth homotopy from the identity to the constant point
\[
p_0=[(e_1,1_G)].
\]

Define the shift operator
\[
S(x_1,x_2,x_3,\ldots)=(0,x_1,x_2,\ldots).
\]
On labelled points, set
\[
S[(x_1,g_1),(x_2,g_2),\ldots]
=
[(0,1_G),(x_1,g_1),(x_2,g_2),\ldots].
\]
This is well-defined because the first inserted coordinate has coefficient zero, hence its group label is irrelevant.

We first construct a homotopy from the identity to the shift. For \(s\in[0,1]\), define
\[
H_s^{(1)}([(x_i,g_i)])
=
\left[
\frac{(1-s)x+sSx}{\|(1-s)x+sSx\|},
\ \text{with transported labels}
\right].
\]
The denominator never vanishes. Indeed, \(x\) and \(Sx\) are orthogonal in the Hilbert space of finitely supported real sequences, and \(\|x\|=\|Sx\|=1\). Hence
\[
\|(1-s)x+sSx\|^2
=
(1-s)^2+s^2,
\]
which is strictly positive for all \(s\in[0,1]\).

The expression has finite support whenever \(x\) has finite support. On each finite chart \(S_I\times G^I\), the formula is a smooth map into a larger finite chart. Hence \(H^{(1)}\) is smooth for the final diffeology.

At \(s=0\), \(H_s^{(1)}\) is the identity. At \(s=1\), it is the shift.

We now contract the shift to \(p_0\). Define
\[
H_s^{(2)}([(x_i,g_i)])
=
\left[
\frac{(1-s)Sx+s e_1}{\|(1-s)Sx+s e_1\|},
\ \text{with transported labels}
\right].
\]
Again the denominator never vanishes. Since \(Sx\) has first coordinate zero, it is orthogonal to \(e_1\), and therefore
\[
\|(1-s)Sx+s e_1\|^2
=
(1-s)^2+s^2.
\]
At \(s=0\) we obtain the shift, and at \(s=1\) we obtain the constant point \(p_0\).

As before, the homotopy is smooth on finite charts, hence smooth for the final diffeology. Concatenating \(H^{(1)}\) and \(H^{(2)}\) gives a smooth contraction of \(E^{\mathrm{sph}}(G)\).
\end{proof}

\subsection{The universal projective class}

\begin{proposition}
The projection
\[
q:E^{\mathrm{sph}}(G)\longrightarrow E^{\mathrm{proj}}(G)
\]
is a principal \(\mathbb Z_2\)-bundle in the diffeological sense.
\end{proposition}

\begin{proof}
The map \(q\) is the quotient map for a smooth free action of the discrete diffeological group \(\mathbb Z_2\). The action is free by the preceding proposition.

It remains to verify the principal bundle condition. Consider the map
\[
E^{\mathrm{sph}}(G)\times \mathbb Z_2
\longrightarrow
E^{\mathrm{sph}}(G)\times_{E^{\mathrm{proj}}(G)}E^{\mathrm{sph}}(G),
\qquad
(p,\varepsilon)\longmapsto (p,\varepsilon p).
\]
It is smooth because the action is smooth. It is bijective because two points lie in the same fiber of \(q\) precisely when they differ by a unique element of \(\mathbb Z_2\), since the action is free. Its inverse sends \((p,p')\) to the unique \(\varepsilon\in\mathbb Z_2\) satisfying \(p'=\varepsilon p\). Since \(\mathbb Z_2\) is discrete and the quotient diffeology is final, this inverse is smooth. Therefore \(q\) is a principal \(\mathbb Z_2\)-bundle.
\end{proof}

\begin{theorem}
The quotient map
\[
\pi:
E^{\mathrm{sph}}(G)
\longrightarrow
E^{\mathrm{proj}}(G)
=
E^{\mathrm{sph}}(G)/\mathbb Z_2
\]
defines a principal \(\mathbb Z_2\)-bundle in the diffeological sense.

Consequently, it determines a canonical characteristic class
\[
\alpha_{\mathrm{sph}}
\in
H^1(E^{\mathrm{proj}}(G);\mathbb Z_2).
\]
\end{theorem}

\begin{proof}
Recall that \(\mathbb Z_2=\{\pm1\}\) acts on
\(E^{\mathrm{sph}}(G)\) by
\[
\varepsilon\cdot(x_i,g_i)
=
(\varepsilon x_i,g_i).
\]

We first show that this action is smooth, free, and proper in the
diffeological sense.

\medskip

\noindent
\textbf{Smoothness.}
On each generating chart
\[
S_I\times G^I,
\]
the action is induced by
\[
\varepsilon\cdot((x_i)_{i\in I},(g_i)_{i\in I})
=
((\varepsilon x_i)_{i\in I},(g_i)_{i\in I}),
\]
which is smooth since multiplication by \(\pm1\) is smooth on the Euclidean
space \(\mathbb R^I\). By the definition of the final diffeology, the induced
action on \(E^{\mathrm{sph}}(G)\) is smooth.

\medskip

\noindent
\textbf{Freeness.}
Assume that
\[
\varepsilon\cdot(x_i,g_i)=(x_i,g_i)
\]
in \(E^{\mathrm{sph}}(G)\). By definition of the quotient relation, this means
that for every index \(i\) such that \(x_i\neq0\),
\[
\varepsilon x_i=x_i.
\]
Since
\[
\sum_i x_i^2=1,
\]
at least one coordinate satisfies \(x_i\neq0\). Hence
\[
\varepsilon=1.
\]
Therefore the action is free.

\medskip

\noindent
\textbf{Quotient structure.}
By definition,
\[
E^{\mathrm{proj}}(G)
=
E^{\mathrm{sph}}(G)/\mathbb Z_2
\]
is equipped with the quotient diffeology induced by the projection
\[
\pi:E^{\mathrm{sph}}(G)\to E^{\mathrm{proj}}(G).
\]

Since the action is smooth and free, the quotient map defines a principal
\(\mathbb Z_2\)-bundle in the diffeological sense.

\medskip

We now construct the characteristic class.

Choose a \(D\)-open cover
\[
(U_a)_a
\]
of \(E^{\mathrm{proj}}(G)\) over which the principal bundle admits smooth local
sections
\[
s_a:U_a\to E^{\mathrm{sph}}(G).
\]

On overlaps
\[
U_a\cap U_b,
\]
there exists a unique locally constant function
\[
g_{ab}:U_a\cap U_b\to\mathbb Z_2
\]
such that
\[
s_b=g_{ab}\cdot s_a.
\]

Since \(\mathbb Z_2\) is discrete, the family \((g_{ab})\) defines a
Čech \(1\)-cocycle with values in \(\mathbb Z_2\). Its cohomology class
\[
[(g_{ab})]
\in
\check H^1(E^{\mathrm{proj}}(G);\mathbb Z_2)
\]
depends only on the isomorphism class of the principal bundle.

Under the standard identification between principal \(\mathbb Z_2\)-bundles
and degree-one cohomology classes, this cocycle defines the characteristic
class
\[
\alpha_{\mathrm{sph}}
\in
H^1(E^{\mathrm{proj}}(G);\mathbb Z_2).
\]
\end{proof}

\begin{remark}
The class
\[
\alpha_{\mathrm{sph}}
\]
is the analogue of the universal first Stiefel--Whitney class associated with
the universal double covering
\[
S^\infty\to\mathbb RP^\infty.
\]
\end{remark}

\begin{remark}
If \(E^{\mathrm{sph}}(G)\) is diffeologically contractible, then
\(E^{\mathrm{proj}}(G)\) plays the role of a diffeological classifying space
for principal \(\mathbb Z_2\)-bundles.
\end{remark}

\begin{remark}
The corresponding cohomology class may also be interpreted as the obstruction
to lifting maps
\[
X\to E^{\mathrm{proj}}(G)
\]
through the double covering
\[
E^{\mathrm{sph}}(G)\to E^{\mathrm{proj}}(G).
\]
\end{remark}

\subsection{The \(G\)-action and the product quotient}

The group \(G\) acts on \(E^{\mathrm{sph}}(G)\) by
\[
h\cdot [(x_i,g_i)]
=
[(x_i,hg_i)].
\]

\begin{proposition}
The actions of \(G\) and \(\mathbb Z_2\) on \(E^{\mathrm{sph}}(G)\) commute.
\end{proposition}

\begin{proof}
For \(h\in G\) and \(\varepsilon\in\mathbb Z_2\),
\[
h\cdot(\varepsilon\cdot[(x_i,g_i)])
=
h\cdot[(\varepsilon x_i,g_i)]
=
[(\varepsilon x_i,hg_i)].
\]
On the other hand,
\[
\varepsilon\cdot(h\cdot[(x_i,g_i)])
=
\varepsilon\cdot[(x_i,hg_i)]
=
[(\varepsilon x_i,hg_i)].
\]
The two expressions agree.
\end{proof}

\begin{definition}
The product spherical classifying quotient is
\[
B^{\mathrm{sph}}_{\times}(G)
:=
E^{\mathrm{sph}}(G)/(G\times \mathbb Z_2).
\]
\end{definition}

\begin{proposition}
There are canonical diffeomorphisms
\[
(E^{\mathrm{sph}}(G)/G)/\mathbb Z_2
\cong
(E^{\mathrm{sph}}(G)/\mathbb Z_2)/G
\cong
B^{\mathrm{sph}}_{\times}(G).
\]
\end{proposition}

\begin{proof}
Since the two actions commute, they define a smooth action of the product group
\[
G\times\mathbb Z_2.
\]
The orbit of a point under first quotienting by \(G\) and then by \(\mathbb Z_2\) is exactly its orbit under \(G\times\mathbb Z_2\). The same is true if the order of the quotients is reversed.

The set-theoretic identifications are therefore clear. The diffeological identifications follow from the universal property of quotient diffeologies: a map out of any of these quotients is smooth if and only if its pullback to \(E^{\mathrm{sph}}(G)\) is smooth and invariant under both actions. This condition is independent of the order in which the two quotients are taken.
\end{proof}

\subsection{Twisted quotients and semidirect products}

The preceding construction corresponds to the product group \(G\times\mathbb Z_2\). To obtain genuinely twisted \(G\)-bundles, one must introduce an action of \(\mathbb Z_2\) on \(G\).

Let
\[
\sigma:\mathbb Z_2\longrightarrow \operatorname{Aut}(G)
\]
be a smooth action. Write
\[
G\rtimes_\sigma \mathbb Z_2
\]
for the corresponding semidirect product.

\begin{definition}
The \(\sigma\)-twisted action of \(G\rtimes_\sigma\mathbb Z_2\) on \(E^{\mathrm{sph}}(G)\) is defined by
\[
(h,\varepsilon)\cdot[(x_i,g_i)]
=
[(\varepsilon x_i, h\,\sigma^\varepsilon(g_i))].
\]
\end{definition}

\begin{proposition}
This defines a smooth action of \(G\rtimes_\sigma\mathbb Z_2\) on \(E^{\mathrm{sph}}(G)\).
\end{proposition}

\begin{proof}
We first check that the formula is compatible with the quotient relation. If
\[
(x_i,g_i)\sim(x_i,g_i'),
\]
then \(g_i=g_i'\) whenever \(x_i\neq0\). Since \(\varepsilon x_i\neq0\) if and only if \(x_i\neq0\), we get
\[
h\,\sigma^\varepsilon(g_i)
=
h\,\sigma^\varepsilon(g_i')
\]
for all active indices. Hence the formula is well-defined.

We now verify the group law. The multiplication in \(G\rtimes_\sigma\mathbb Z_2\) is
\[
(h,\varepsilon)(k,\eta)
=
(h\,\sigma^\varepsilon(k),\varepsilon\eta).
\]
Then
\[
(k,\eta)\cdot[(x_i,g_i)]
=
[(\eta x_i,k\,\sigma^\eta(g_i))].
\]
Applying \((h,\varepsilon)\) gives
\[
(h,\varepsilon)\cdot[(\eta x_i,k\,\sigma^\eta(g_i))]
=
[(\varepsilon\eta x_i,
h\,\sigma^\varepsilon(k\,\sigma^\eta(g_i)))].
\]
Since \(\sigma^\varepsilon\) is an automorphism,
\[
\sigma^\varepsilon(k\,\sigma^\eta(g_i))
=
\sigma^\varepsilon(k)\,\sigma^{\varepsilon\eta}(g_i).
\]
Thus
\[
(h,\varepsilon)\cdot((k,\eta)\cdot[(x_i,g_i)])
=
[(\varepsilon\eta x_i,
h\,\sigma^\varepsilon(k)\,\sigma^{\varepsilon\eta}(g_i))],
\]
which is precisely
\[
(h\,\sigma^\varepsilon(k),\varepsilon\eta)\cdot[(x_i,g_i)].
\]
Hence this is an action.

Smoothness follows from smoothness of the \(G\)-action, of the finite \(\mathbb Z_2\)-action, and of the automorphism \(\sigma^\varepsilon\) on each finite chart \(S_I\times G^I\).
\end{proof}

\begin{definition}
The \(\sigma\)-twisted spherical classifying quotient is
\[
B^{\mathrm{sph}}_\sigma(G)
:=
E^{\mathrm{sph}}(G)/(G\rtimes_\sigma\mathbb Z_2).
\]
\end{definition}

\begin{remark}
When \(\sigma\) is trivial, \(G\rtimes_\sigma\mathbb Z_2=G\times\mathbb Z_2\), and \(B^{\mathrm{sph}}_\sigma(G)=B^{\mathrm{sph}}_{\times}(G)\).
\end{remark}

\begin{remark}
For nontrivial \(\sigma\), the \(G\)- and \(\mathbb Z_2\)-parts no longer commute as independent actions. Thus the twisted quotient must be understood as a quotient by the semidirect product, not as an iterated quotient by commuting actions.
\end{remark}

\subsection{Cocycle interpretation}

Let \(X\) be a diffeological space and let
\[
P\longrightarrow X
\]
be a principal \(G\rtimes_\sigma\mathbb Z_2\)-bundle. Choose a cover \((U_i)\) and transition functions
\[
h_{ij}:U_{ij}\longrightarrow G\rtimes_\sigma\mathbb Z_2.
\]
Write
\[
h_{ij}=(g_{ij},\varepsilon_{ij}).
\]
The cocycle condition
\[
h_{ij}h_{jk}=h_{ik}
\]
becomes
\[
\varepsilon_{ij}\varepsilon_{jk}=\varepsilon_{ik},
\]
and
\[
g_{ij}\,\sigma^{\varepsilon_{ij}}(g_{jk})=g_{ik}.
\]

\begin{proposition}
The \(\mathbb Z_2\)-components \((\varepsilon_{ij})\) define a class
\[
\alpha(P)\in H^1(X;\mathbb Z_2).
\]
The \(G\)-components \((g_{ij})\) form a \(G\)-valued cocycle twisted by \(\alpha(P)\).
\end{proposition}

\begin{proof}
The equation
\[
\varepsilon_{ij}\varepsilon_{jk}=\varepsilon_{ik}
\]
is exactly the Čech \(1\)-cocycle condition for \(\mathbb Z_2\)-valued transition functions. Hence it defines a class in \(H^1(X;\mathbb Z_2)\).

The second equation
\[
g_{ij}\,\sigma^{\varepsilon_{ij}}(g_{jk})=g_{ik}
\]
says precisely that the \(G\)-valued transition functions fail to satisfy the ordinary cocycle condition by the twisting induced by the \(\mathbb Z_2\)-cocycle. This is the usual cocycle description of a \(G\)-bundle twisted by the local system determined by \(\alpha(P)\).
\end{proof}

\begin{remark}
Thus the topological information carried by the projective spherical quotient is not merely the product of a \(G\)-bundle and a double cover. In the presence of a nontrivial action \(\sigma\), it encodes a genuinely twisted principal \(G\)-bundle.
\end{remark}

\section{Classifying properties, characteristic cocycles and canonical forms}
\label{sec:classifying-properties}

\subsection{Classifying property of the spherical quotient}



\begin{definition}
Let
\[
\sigma:\mathbb Z_2\to \mathrm{Aut}(G)
\]
be a smooth action of \(\mathbb Z_2\) on the diffeological group \(G\).

We denote by
\[
H_\sigma
=
G\rtimes_\sigma \mathbb Z_2
\]
the corresponding semidirect product.

The \emph{\(\sigma\)-twisted spherical classifying space} is defined by
\[
B^{\mathrm{sph}}_\sigma(G)
:=
E^{\mathrm{sph}}(G)/H_\sigma,
\]
where \(H_\sigma\) acts on \(E^{\mathrm{sph}}(G)\) through the combined
\(G\)-action and \(\mathbb Z_2\)-action.
\end{definition}

\begin{theorem}
Assume that:
\begin{itemize}
\item \(E^{\mathrm{sph}}(G)\) is smoothly contractible;
\item the action of \(H_\sigma\) on \(E^{\mathrm{sph}}(G)\) is smooth and free.
\end{itemize}

Then the quotient projection
\[
\pi:
E^{\mathrm{sph}}(G)
\longrightarrow
B^{\mathrm{sph}}_\sigma(G)
\]
is a principal \(H_\sigma\)-bundle in the diffeological sense.

Moreover, it is universal for \(D\)-numerable principal
\(H_\sigma\)-bundles: for every diffeological space \(X\), smooth homotopy
classes of maps
\[
[X,B^{\mathrm{sph}}_\sigma(G)]^\infty
\]
are in bijection with isomorphism classes of \(D\)-numerable principal
\(H_\sigma\)-bundles over \(X\).
\end{theorem}

\begin{proof}
We first prove that
\[
\pi:
E^{\mathrm{sph}}(G)
\to
B^{\mathrm{sph}}_\sigma(G)
\]
is a principal \(H_\sigma\)-bundle.

\medskip

By assumption, the action
\[
H_\sigma\times E^{\mathrm{sph}}(G)
\to
E^{\mathrm{sph}}(G)
\]
is smooth and free. The quotient space
\[
B^{\mathrm{sph}}_\sigma(G)
=
E^{\mathrm{sph}}(G)/H_\sigma
\]
is equipped with the quotient diffeology induced by the projection \(\pi\).

We now consider the canonical map
\[
\Theta:
E^{\mathrm{sph}}(G)\times H_\sigma
\longrightarrow
E^{\mathrm{sph}}(G)\times_{B^{\mathrm{sph}}_\sigma(G)}
E^{\mathrm{sph}}(G),
\]
defined by
\[
\Theta(p,h)
=
(p,p\cdot h).
\]

We show that \(\Theta\) is a diffeomorphism.

\medskip

\noindent
\textbf{Injectivity.}
Suppose
\[
\Theta(p,h_1)=\Theta(p,h_2).
\]
Then
\[
p\cdot h_1=p\cdot h_2.
\]
Since the action is free,
\[
h_1=h_2.
\]
Hence \(\Theta\) is injective.

\medskip

\noindent
\textbf{Surjectivity.}
Let
\[
(p,q)
\in
E^{\mathrm{sph}}(G)\times_{B^{\mathrm{sph}}_\sigma(G)}
E^{\mathrm{sph}}(G).
\]
By definition of the fiber product,
\[
\pi(p)=\pi(q).
\]
Therefore \(p\) and \(q\) belong to the same \(H_\sigma\)-orbit. Hence there
exists
\[
h\in H_\sigma
\]
such that
\[
q=p\cdot h.
\]
Thus
\[
(p,q)=\Theta(p,h),
\]
and \(\Theta\) is surjective.

\medskip

\noindent
\textbf{Smoothness.}
The map \(\Theta\) is smooth because the action is smooth.

Its inverse is also smooth. Indeed, given
\[
(p,q)
\in
E^{\mathrm{sph}}(G)\times_{B^{\mathrm{sph}}_\sigma(G)}
E^{\mathrm{sph}}(G),
\]
the unique element
\[
h\in H_\sigma
\]
satisfying
\[
q=p\cdot h
\]
depends smoothly on \((p,q)\) by the definition of the quotient diffeology
and the freeness of the action.

Therefore \(\Theta\) is a diffeomorphism, and \(\pi\) is a principal
\(H_\sigma\)-bundle in the diffeological sense.

\medskip

We now prove universality.

Since \(E^{\mathrm{sph}}(G)\) is smoothly contractible, the principal bundle
\[
E^{\mathrm{sph}}(G)
\to
B^{\mathrm{sph}}_\sigma(G)
\]
satisfies the defining property of a universal bundle in the diffeological
Milnor framework developed in \cite{MagnotWatts2017}.

More precisely, the diffeological classification theorem states that if
\(E\to B\) is a smoothly contractible principal \(H_\sigma\)-bundle, then for
every diffeological space \(X\), every \(D\)-numerable principal
\(H_\sigma\)-bundle over \(X\) is isomorphic to the pull-back of
\(E\to B\) along a smooth map
\[
f:X\to B,
\]
and two such pull-backs are isomorphic if and only if the corresponding maps
are smoothly homotopic.

Applying this result to
\[
E^{\mathrm{sph}}(G)
\to
B^{\mathrm{sph}}_\sigma(G)
\]
gives the desired classification statement.
\end{proof}

\begin{remark}
When the action \(\sigma\) is trivial, one recovers the untwisted spherical
classifying space associated with
\[
G\times \mathbb Z_2.
\]
\end{remark}

\begin{remark}
The twisted quotient encodes simultaneously:
\begin{itemize}
\item the topology of the group \(G\),
\item the projective \(\mathbb Z_2\)-symmetry,
\item and the interaction between both structures through the twisting
homomorphism \(\sigma\).
\end{itemize}
\end{remark}

\begin{remark}
The associated characteristic classes naturally live in the cohomology of
\[
B^{\mathrm{sph}}_\sigma(G),
\]
and may be interpreted as twisted analogues of Stiefel--Whitney-type
obstructions.
\end{remark}
\subsection{Explicit construction of classifying maps}

Let
\[
P\longrightarrow X
\]
be a \(D\)-numerable principal \(H_\sigma\)-bundle.

Choose:
\begin{itemize}
\item a locally finite open cover \((U_i)_{i\in I}\) of \(X\),
\item local sections \(s_i:U_i\to P\),
\item a smooth partition of unity \((\lambda_i)_{i\in I}\) subordinate to \((U_i)\).
\end{itemize}

Let the transition functions be
\[
h_{ij}:U_i\cap U_j\to H_\sigma,
\qquad
s_j=s_i h_{ij}.
\]
Write
\[
h_{ij}=(g_{ij},\varepsilon_{ij})
\]
with
\[
g_{ij}:U_{ij}\to G,
\qquad
\varepsilon_{ij}:U_{ij}\to \mathbb Z_2.
\]

\begin{definition}
The classifying map
\[
f_P:X\longrightarrow B^{\mathrm{sph}}_\sigma(G)
\]
is defined locally by
\[
f_P(x)
=
\left[
\left(\sqrt{\lambda_i(x)},h_i(x)\right)_{i\in I}
\right],
\]
where \(h_i(x)\in H_\sigma\) are the local \(H_\sigma\)-coordinates determined by the chosen trivializations.
\end{definition}

Equivalently, in a local trivialization over \(U_k\), one may write
\[
f_P(x)
=
\left[
\left(\sqrt{\lambda_i(x)},h_{ki}(x)\right)_{i\in I}
\right].
\]

\begin{proposition}
The map \(f_P\) is well-defined and smooth.
\end{proposition}

\begin{proof}
On \(U_k\), the formula
\[
x\longmapsto
\left(\sqrt{\lambda_i(x)},h_{ki}(x)\right)_{i\in I}
\]
defines a map into a finite spherical chart because the partition of unity is locally finite and
\[
\sum_i \lambda_i(x)=1.
\]
Thus
\[
\sum_i \bigl(\sqrt{\lambda_i(x)}\bigr)^2=1.
\]

The square roots are smooth because the functions \(\lambda_i\) are chosen in the \(D\)-numerable sense, hence are stationary at the boundary of their supports in the diffeological setting. Therefore the map locally defines a plot of \(E^{\mathrm{sph}}(G)\), and its projection to the quotient gives a smooth map into \(B^{\mathrm{sph}}_\sigma(G)\).

On overlaps \(U_k\cap U_\ell\), the two expressions differ by the transition function \(h_{\ell k}\in H_\sigma\). Since \(B^{\mathrm{sph}}_\sigma(G)\) is the quotient by \(H_\sigma\), they define the same point. Hence the local formulas glue to a globally defined smooth map.
\end{proof}

\begin{proposition}
The pull-back bundle
\[
f_P^*E^{\mathrm{sph}}(G)
\]
is isomorphic to \(P\) as a principal \(H_\sigma\)-bundle.
\end{proposition}

\begin{proof}
The pull-back is described locally by the same transition functions as the universal bundle. In the chart \(U_k\), the universal point is represented by
\[
\left(\sqrt{\lambda_i(x)},h_{ki}(x)\right)_{i\in I}.
\]
On \(U_k\cap U_\ell\), the change of representative is precisely multiplication by
\[
h_{\ell k}(x).
\]
Therefore the transition functions of \(f_P^*E^{\mathrm{sph}}(G)\) coincide with the transition functions of \(P\). Hence the two principal bundles are isomorphic.
\end{proof}

\subsection{The associated \texorpdfstring{\(\mathbb Z_2\)}{Z2}-class}

The projection
\[
H_\sigma=G\rtimes_\sigma\mathbb Z_2
\longrightarrow
\mathbb Z_2
\]
induces a map of classifying spaces
\[
B^{\mathrm{sph}}_\sigma(G)
\longrightarrow
B\mathbb Z_2.
\]

\begin{definition}
Let
\[
\alpha_{\mathrm{univ}}\in H^1(B^{\mathrm{sph}}_\sigma(G);\mathbb Z_2)
\]
be the pull-back of the universal class in
\[
H^1(B\mathbb Z_2;\mathbb Z_2).
\]
\end{definition}

\begin{proposition}
For a principal \(H_\sigma\)-bundle \(P\to X\) classified by \(f_P\), the class
\[
f_P^*\alpha_{\mathrm{univ}}\in H^1(X;\mathbb Z_2)
\]
is represented by the Čech cocycle \((\varepsilon_{ij})\).
\end{proposition}

\begin{proof}
The \(\mathbb Z_2\)-bundle associated with \(P\) is obtained by quotienting out the \(G\)-part:
\[
P/G\longrightarrow X.
\]
Its transition functions are precisely the \(\mathbb Z_2\)-components \(\varepsilon_{ij}\) of the \(H_\sigma\)-transition functions
\[
h_{ij}=(g_{ij},\varepsilon_{ij}).
\]
Thus \((\varepsilon_{ij})\) represents the first cohomology class associated with the induced double cover. Since \(f_P\) classifies \(P\), this class is the pull-back of the universal class.
\end{proof}

\begin{remark}
Thus the class \(f_P^*\alpha_{\mathrm{univ}}\) is the topological obstruction measuring the failure of the \(H_\sigma\)-bundle to reduce to an ordinary \(G\)-bundle.
\end{remark}

\subsection{Twisted cocycles}

The cocycle condition for \(h_{ij}\) reads
\[
h_{ij}h_{jk}=h_{ik}.
\]
Writing
\[
h_{ij}=(g_{ij},\varepsilon_{ij}),
\]
one obtains
\[
\varepsilon_{ij}\varepsilon_{jk}=\varepsilon_{ik},
\]
and
\[
g_{ij}\sigma^{\varepsilon_{ij}}(g_{jk})=g_{ik}.
\]

\begin{proposition}
Let
\[
H_\sigma=G\rtimes_\sigma \mathbb Z_2
\]
with product
\[
(g,\varepsilon)(g',\varepsilon')
=
(g\,\sigma_\varepsilon(g'),\varepsilon\varepsilon').
\]
Let \(P\to X\) be a principal \(H_\sigma\)-bundle, trivialized over an open
cover \((U_i)_i\). Write its transition functions as
\[
h_{ij}=(g_{ij},\varepsilon_{ij})
:
U_{ij}\to G\rtimes_\sigma\mathbb Z_2.
\]
Then the functions \(\varepsilon_{ij}\) define a \(\mathbb Z_2\)-valued
Čech cocycle, hence a class
\[
\alpha(P)=[\varepsilon_{ij}]\in H^1(X;\mathbb Z_2),
\]
and the functions \(g_{ij}\) satisfy the twisted cocycle relation
\[
g_{ij}\,\sigma_{\varepsilon_{ij}}(g_{jk})
=
g_{ik}
\qquad
\text{on }U_{ijk}.
\]
Thus the pair
\[
(\varepsilon_{ij},g_{ij})
\]
defines a non-abelian \(G\)-valued cocycle twisted by the local system
associated with \(\alpha(P)\).
\end{proposition}

\begin{proof}
Since \(P\to X\) is a principal \(H_\sigma\)-bundle, its transition functions
satisfy the usual non-abelian Čech cocycle condition
\[
h_{ij}h_{jk}=h_{ik}
\qquad
\text{on }U_{ijk}.
\]
Writing
\[
h_{ij}=(g_{ij},\varepsilon_{ij}),
\qquad
h_{jk}=(g_{jk},\varepsilon_{jk}),
\]
and using the product in the semidirect product, we obtain
\[
h_{ij}h_{jk}
=
(g_{ij},\varepsilon_{ij})(g_{jk},\varepsilon_{jk})
=
(g_{ij}\sigma_{\varepsilon_{ij}}(g_{jk}),
\varepsilon_{ij}\varepsilon_{jk}).
\]
Therefore the equality
\[
h_{ij}h_{jk}=h_{ik}
=
(g_{ik},\varepsilon_{ik})
\]
is equivalent to the two equations
\[
\varepsilon_{ij}\varepsilon_{jk}
=
\varepsilon_{ik},
\]
and
\[
g_{ij}\sigma_{\varepsilon_{ij}}(g_{jk})
=
g_{ik}.
\]

The first equation is exactly the Čech cocycle condition for the
\(\mathbb Z_2\)-valued transition functions \((\varepsilon_{ij})\). Hence it
defines a class
\[
\alpha(P)\in H^1(X;\mathbb Z_2).
\]

The second equation says that the \(G\)-valued transition functions do not
satisfy the ordinary cocycle condition unless the action of \(\mathbb Z_2\)
is trivial. Instead, the term \(g_{jk}\) is transported by the automorphism
\[
\sigma_{\varepsilon_{ij}}\in \operatorname{Aut}(G).
\]
Thus the \(G\)-valued part is a non-abelian cocycle with coefficients twisted
by the \(\mathbb Z_2\)-local system represented by
\[
(\varepsilon_{ij}).
\]
This proves the claim.
\end{proof}
\begin{remark}
If \(\sigma\) is trivial, the relation reduces to the ordinary non-abelian
cocycle condition
\[
g_{ij}g_{jk}=g_{ik}.
\]
Thus the above construction is exactly the usual Čech description of a
principal \(G\)-bundle, modified by the \(\mathbb Z_2\)-twisting.
\end{remark}
\subsection{Pull-back of the canonical connection}

Let
\[
\Theta
\]
denote the canonical connection form on the universal spherical bundle
\[
E^{\mathrm{sph}}(G)\to B^{\mathrm{sph}}_\sigma(G).
\]

Locally, in the chart indexed by \(I\), this connection is given by the barycentric formula
\[
\Theta
=
\sum_{i\in I} x_i^2\,\theta_i,
\]
where \(\theta_i\) is the Maurer--Cartan form on the \(i\)-th copy of \(G\) or, more generally, on the \(i\)-th copy of \(H_\sigma\).

\begin{proposition}
Let \(P\to X\) be classified by \(f_P:X\to B^{\mathrm{sph}}_\sigma(G)\). Then the pull-back connection
\[
\nabla_P:=f_P^*\Theta
\]
is locally given by
\[
A_k
=
\sum_i \lambda_i\, h_{ki}^{-1}dh_{ki}
\]
on \(U_k\).
\end{proposition}

\begin{proof}
On \(U_k\), the classifying map is represented by
\[
x\longmapsto
\left(\sqrt{\lambda_i(x)},h_{ki}(x)\right)_{i\in I}.
\]
Pulling back the universal connection
\[
\Theta=\sum_i x_i^2\theta_i
\]
gives
\[
f_P^*\Theta
=
\sum_i \lambda_i\, h_{ki}^*\theta,
\]
where \(\theta\) is the Maurer--Cartan form. Since
\[
h_{ki}^*\theta=h_{ki}^{-1}dh_{ki},
\]
we obtain
\[
A_k=\sum_i \lambda_i\,h_{ki}^{-1}dh_{ki}.
\]
\end{proof}

\begin{remark}
This is the spherical analogue of the barycentric connection in the Milnor--Watts construction. The only difference is the replacement
\[
t_i=x_i^2.
\]
\end{remark}

\subsection{Curvature and characteristic forms}

Let
\[
A_k
=
\sum_i \lambda_i\,h_{ki}^{-1}dh_{ki}
\]
be the local connection form. Its curvature is
\[
F_k
=
dA_k+\frac12[A_k,A_k].
\]

\begin{proposition}
The local curvature forms \(F_k\) transform under the usual gauge transformation rule on overlaps:
\[
F_\ell
=
\operatorname{Ad}_{h_{\ell k}^{-1}}F_k.
\]
\end{proposition}

\begin{proof}
On overlaps, the local connection forms satisfy
\[
A_\ell
=
\operatorname{Ad}_{h_{\ell k}^{-1}}A_k
+
h_{\ell k}^{-1}dh_{\ell k}.
\]
This follows from the transformation law of the Maurer--Cartan form and the cocycle identity for \(h_{ij}\). Taking curvature gives the standard transformation law
\[
F_\ell
=
\operatorname{Ad}_{h_{\ell k}^{-1}}F_k.
\]
\end{proof}

Let
\[
P\in \operatorname{Sym}^r(\mathfrak h_\sigma^*)^{H_\sigma}
\]
be an invariant polynomial on
\[
\mathfrak h_\sigma=\operatorname{Lie}(H_\sigma).
\]
Since \(\mathbb Z_2\) is discrete,
\[
\mathfrak h_\sigma\simeq \mathfrak g.
\]

\begin{definition}
The characteristic form associated with \(P\) is
\[
\omega_P(F)
=
P(F_k,\ldots,F_k)
\]
on \(U_k\).
\end{definition}

\begin{theorem}
The local forms
\[
P(F_k,\ldots,F_k)
\]
glue to a global closed differential form on \(X\). Its de Rham cohomology class is independent of the choice of the local trivializations and of the representative connection.
\end{theorem}

\begin{proof}
By the curvature transformation rule and the invariance of \(P\), one has
\[
P(F_\ell,\ldots,F_\ell)
=
P(F_k,\ldots,F_k)
\]
on overlaps. Thus the local forms glue to a global form.

The Bianchi identity gives
\[
d_A F=0.
\]
By the standard Chern--Weil argument, this implies
\[
d\,P(F,\ldots,F)=0.
\]
If the connection is changed through a smooth one-parameter family, the derivative of the corresponding characteristic form is exact by the transgression formula. Hence the de Rham cohomology class is independent of the connection.
\end{proof}

\begin{remark}
In the twisted case, the invariant polynomial must be invariant not only under the adjoint action of \(G\), but also under the induced action of \(\mathbb Z_2\) on \(\mathfrak g\).
\end{remark}

\subsection{Relation with known topological invariants}

The projection
\[
H_\sigma\to \mathbb Z_2
\]
gives the degree-one invariant
\[
\alpha(P)\in H^1(X;\mathbb Z_2).
\]

The inclusion of \(G\) into \(H_\sigma\) and the invariant polynomials on \(\mathfrak g\) yield the usual characteristic classes of the underlying \(G\)-geometry, whenever they are invariant under the \(\mathbb Z_2\)-action.

\begin{theorem}
Let \(P\to X\) be a principal \(H_\sigma\)-bundle classified by
\[
f_P:X\to B^{\mathrm{sph}}_\sigma(G).
\]
Then:
\begin{enumerate}
\item the class \(f_P^*\alpha_{\mathrm{univ}}\) is the first \(\mathbb Z_2\)-obstruction class of \(P\);
\item the characteristic forms obtained by pulling back the universal barycentric connection coincide with the Chern--Weil forms of the associated \(H_\sigma\)-connection;
\item when \(\sigma\) is trivial, these invariants split into the ordinary characteristic classes of the \(G\)-bundle together with the independent class in \(H^1(X;\mathbb Z_2)\).
\end{enumerate}
\end{theorem}

\begin{proof}
The first point follows from the preceding identification of \(f_P^*\alpha_{\mathrm{univ}}\) with the Čech class of the \(\mathbb Z_2\)-components \((\varepsilon_{ij})\).

The second point follows from the explicit pull-back formula
\[
A_k=\sum_i\lambda_i h_{ki}^{-1}dh_{ki},
\]
which is precisely the local expression for the induced connection on \(P\). The characteristic forms are therefore the usual Chern--Weil forms of this connection.

If \(\sigma\) is trivial, then
\[
H_\sigma=G\times \mathbb Z_2.
\]
Thus every principal \(H_\sigma\)-bundle is equivalent to a pair consisting of a principal \(G\)-bundle and a principal \(\mathbb Z_2\)-bundle. The corresponding invariants split accordingly.
\end{proof}

\section{Gerbes and secondary obstruction classes}
\label{sec:gerbes}

\subsection{Central extensions compatible with the twist}

Let
\[
1\longrightarrow A \longrightarrow \widehat G
\stackrel{q}{\longrightarrow} G \longrightarrow 1
\]
be a central extension of diffeological groups, with \(A\) abelian.

Assume that the action
\[
\sigma:\mathbb Z_2\longrightarrow \operatorname{Aut}(G)
\]
lifts to an action
\[
\widehat\sigma:\mathbb Z_2\longrightarrow \operatorname{Aut}(\widehat G)
\]
such that
\[
q\circ \widehat\sigma^\varepsilon
=
\sigma^\varepsilon\circ q
\]
for every \(\varepsilon\in\mathbb Z_2\).

Then we have a central extension
\[
1\longrightarrow A
\longrightarrow
\widehat G\rtimes_{\widehat\sigma}\mathbb Z_2
\longrightarrow
G\rtimes_\sigma \mathbb Z_2
\longrightarrow 1.
\]

Set
\[
H_\sigma:=G\rtimes_\sigma\mathbb Z_2,
\qquad
\widehat H_{\widehat\sigma}:=\widehat G\rtimes_{\widehat\sigma}\mathbb Z_2.
\]

\subsection{The lifting problem}

Let \(P\to X\) be a principal \(H_\sigma\)-bundle over a diffeological space \(X\).

Choose a good cover \((U_i)\) and transition functions
\[
h_{ij}:U_{ij}\to H_\sigma.
\]
Write
\[
h_{ij}=(g_{ij},\varepsilon_{ij}),
\]
where
\[
g_{ij}:U_{ij}\to G,
\qquad
\varepsilon_{ij}:U_{ij}\to\mathbb Z_2.
\]

The cocycle condition is
\[
h_{ij}h_{jk}=h_{ik}.
\]
Equivalently,
\[
\varepsilon_{ij}\varepsilon_{jk}=\varepsilon_{ik},
\]
and
\[
g_{ij}\,\sigma^{\varepsilon_{ij}}(g_{jk})=g_{ik}.
\]

\begin{definition}
A lift of \(P\) to \(\widehat H_{\widehat\sigma}\) is a principal
\(\widehat H_{\widehat\sigma}\)-bundle \(\widehat P\to X\) together with an isomorphism
\[
\widehat P/A \simeq P.
\]
\end{definition}

The obstruction to such a lift is measured by a gerbe.

\subsection{The obstruction cocycle}

Choose local lifts
\[
\widehat g_{ij}:U_{ij}\to \widehat G
\]
of \(g_{ij}\), i.e.
\[
q(\widehat g_{ij})=g_{ij}.
\]

Define
\[
a_{ijk}
:=
\widehat g_{ij}\,
\widehat\sigma^{\varepsilon_{ij}}(\widehat g_{jk})\,
\widehat g_{ik}^{-1}.
\]

\begin{proposition}
For every triple intersection \(U_{ijk}\), the element \(a_{ijk}\) takes values in \(A\).
\end{proposition}

\begin{proof}
Apply \(q:\widehat G\to G\). We get
\[
q(a_{ijk})
=
q(\widehat g_{ij})\,
q\bigl(\widehat\sigma^{\varepsilon_{ij}}(\widehat g_{jk})\bigr)\,
q(\widehat g_{ik}^{-1}).
\]
Using compatibility of \(q\) with the lifted action,
\[
q\bigl(\widehat\sigma^{\varepsilon_{ij}}(\widehat g_{jk})\bigr)
=
\sigma^{\varepsilon_{ij}}(q(\widehat g_{jk}))
=
\sigma^{\varepsilon_{ij}}(g_{jk}).
\]
Thus
\[
q(a_{ijk})
=
g_{ij}\,\sigma^{\varepsilon_{ij}}(g_{jk})\,g_{ik}^{-1}.
\]
By the cocycle condition for \(H_\sigma\),
\[
g_{ij}\,\sigma^{\varepsilon_{ij}}(g_{jk})=g_{ik}.
\]
Therefore
\[
q(a_{ijk})=e.
\]
Hence
\[
a_{ijk}\in \ker(q)=A.
\]
\end{proof}

\subsection{The twisted cocycle identity}

The \(\mathbb Z_2\)-cocycle \((\varepsilon_{ij})\) acts on \(A\) through the lifted action \(\widehat\sigma\). Thus \(A\) becomes a local system, denoted
\[
A_\alpha,
\]
where
\[
\alpha=[\varepsilon_{ij}]\in H^1(X;\mathbb Z_2).
\]

\begin{theorem}
The family \((a_{ijk})\) defines a Čech \(2\)-cocycle with coefficients in the local system \(A_\alpha\).
\end{theorem}

\begin{proof}
We must prove the twisted cocycle identity on quadruple intersections.

By definition,
\[
a_{ijk}
=
\widehat g_{ij}\,
\widehat\sigma^{\varepsilon_{ij}}(\widehat g_{jk})\,
\widehat g_{ik}^{-1}.
\]

Equivalently,
\[
\widehat g_{ij}\,
\widehat\sigma^{\varepsilon_{ij}}(\widehat g_{jk})
=
a_{ijk}\,\widehat g_{ik}.
\]

Now consider the product
\[
\widehat g_{ij}\,
\widehat\sigma^{\varepsilon_{ij}}(\widehat g_{jk})\,
\widehat\sigma^{\varepsilon_{ij}\varepsilon_{jk}}(\widehat g_{kl}).
\]

We compute it in two ways.

First, group the first two factors:
\[
\widehat g_{ij}\,
\widehat\sigma^{\varepsilon_{ij}}(\widehat g_{jk})
=
a_{ijk}\widehat g_{ik}.
\]
Hence
\[
\widehat g_{ij}\,
\widehat\sigma^{\varepsilon_{ij}}(\widehat g_{jk})\,
\widehat\sigma^{\varepsilon_{ij}\varepsilon_{jk}}(\widehat g_{kl})
=
a_{ijk}\,
\widehat g_{ik}\,
\widehat\sigma^{\varepsilon_{ik}}(\widehat g_{kl}).
\]
Using the definition of \(a_{ikl}\),
\[
\widehat g_{ik}\,
\widehat\sigma^{\varepsilon_{ik}}(\widehat g_{kl})
=
a_{ikl}\widehat g_{il}.
\]
Therefore the product equals
\[
a_{ijk}a_{ikl}\widehat g_{il}.
\]

Second, group the last two factors:
\[
\widehat g_{jk}\,
\widehat\sigma^{\varepsilon_{jk}}(\widehat g_{kl})
=
a_{jkl}\widehat g_{jl}.
\]
Applying \(\widehat\sigma^{\varepsilon_{ij}}\), we obtain
\[
\widehat\sigma^{\varepsilon_{ij}}(\widehat g_{jk})\,
\widehat\sigma^{\varepsilon_{ij}\varepsilon_{jk}}(\widehat g_{kl})
=
\widehat\sigma^{\varepsilon_{ij}}(a_{jkl})\,
\widehat\sigma^{\varepsilon_{ij}}(\widehat g_{jl}).
\]
Thus the same product equals
\[
\widehat g_{ij}\,
\widehat\sigma^{\varepsilon_{ij}}(a_{jkl})\,
\widehat\sigma^{\varepsilon_{ij}}(\widehat g_{jl}).
\]
Since \(A\) is central in \(\widehat G\), the term
\(\widehat\sigma^{\varepsilon_{ij}}(a_{jkl})\in A\) can be moved to the front:
\[
=
\widehat\sigma^{\varepsilon_{ij}}(a_{jkl})\,
\widehat g_{ij}\,
\widehat\sigma^{\varepsilon_{ij}}(\widehat g_{jl}).
\]
Using the definition of \(a_{ijl}\),
\[
\widehat g_{ij}\,
\widehat\sigma^{\varepsilon_{ij}}(\widehat g_{jl})
=
a_{ijl}\widehat g_{il}.
\]
Therefore the product equals
\[
\widehat\sigma^{\varepsilon_{ij}}(a_{jkl})\,a_{ijl}\widehat g_{il}.
\]

Comparing the two expressions, we obtain
\[
a_{ijk}a_{ikl}
=
\widehat\sigma^{\varepsilon_{ij}}(a_{jkl})\,a_{ijl}.
\]
This is exactly the Čech \(2\)-cocycle condition with coefficients twisted by the local system \(A_\alpha\).
\end{proof}

\subsection{Independence of choices}

\begin{theorem}
The cohomology class
\[
[a]\in H^2(X;A_\alpha)
\]
defined by \((a_{ijk})\) is independent of the choice of local lifts \(\widehat g_{ij}\).
\end{theorem}

\begin{proof}
Let \(\widehat g_{ij}'\) be another choice of lifts. Since both \(\widehat g_{ij}\) and \(\widehat g_{ij}'\) lift \(g_{ij}\), there exists a map
\[
b_{ij}:U_{ij}\to A
\]
such that
\[
\widehat g_{ij}'=b_{ij}\widehat g_{ij}.
\]

Compute the new cocycle:
\[
a'_{ijk}
=
\widehat g'_{ij}\,
\widehat\sigma^{\varepsilon_{ij}}(\widehat g'_{jk})\,
(\widehat g'_{ik})^{-1}.
\]
Substituting,
\[
a'_{ijk}
=
b_{ij}\widehat g_{ij}\,
\widehat\sigma^{\varepsilon_{ij}}(b_{jk}\widehat g_{jk})\,
(b_{ik}\widehat g_{ik})^{-1}.
\]
Since \(\widehat\sigma^{\varepsilon_{ij}}\) is an automorphism,
\[
\widehat\sigma^{\varepsilon_{ij}}(b_{jk}\widehat g_{jk})
=
\widehat\sigma^{\varepsilon_{ij}}(b_{jk})\,
\widehat\sigma^{\varepsilon_{ij}}(\widehat g_{jk}).
\]
Since \(A\) is central,
\[
a'_{ijk}
=
b_{ij}\,
\widehat\sigma^{\varepsilon_{ij}}(b_{jk})\,
b_{ik}^{-1}
\,
\widehat g_{ij}\widehat\sigma^{\varepsilon_{ij}}(\widehat g_{jk})\widehat g_{ik}^{-1}.
\]
Thus
\[
a'_{ijk}
=
b_{ij}\,
\widehat\sigma^{\varepsilon_{ij}}(b_{jk})\,
b_{ik}^{-1}
\,a_{ijk}.
\]
This is precisely multiplication by a Čech coboundary in the twisted local system \(A_\alpha\). Hence \(a\) and \(a'\) define the same class in \(H^2(X;A_\alpha)\).
\end{proof}

\subsection{Vanishing criterion}

\begin{theorem}
The obstruction class
\[
[a]\in H^2(X;A_\alpha)
\]
vanishes if and only if the principal \(H_\sigma\)-bundle \(P\to X\) lifts to a principal \(\widehat H_{\widehat\sigma}\)-bundle.
\end{theorem}

\begin{proof}
Assume first that \(P\) lifts to a principal \(\widehat H_{\widehat\sigma}\)-bundle. Then there exist transition functions
\[
\widehat h_{ij}=(\widehat g_{ij},\varepsilon_{ij})
\]
satisfying the strict cocycle condition
\[
\widehat h_{ij}\widehat h_{jk}=\widehat h_{ik}.
\]
In terms of the \(\widehat G\)-components, this gives
\[
\widehat g_{ij}\,
\widehat\sigma^{\varepsilon_{ij}}(\widehat g_{jk})
=
\widehat g_{ik}.
\]
Hence
\[
a_{ijk}=1
\]
for all \(i,j,k\), so \([a]=0\).

Conversely, assume \([a]=0\). Then there exist maps
\[
b_{ij}:U_{ij}\to A
\]
such that
\[
a_{ijk}
=
b_{ij}^{-1}\,
\widehat\sigma^{\varepsilon_{ij}}(b_{jk})^{-1}\,
b_{ik}.
\]
Define new lifts
\[
\widehat g'_{ij}:=b_{ij}\widehat g_{ij}.
\]
By the computation in the proof of independence of choices, the new obstruction cocycle is
\[
a'_{ijk}
=
b_{ij}\,
\widehat\sigma^{\varepsilon_{ij}}(b_{jk})\,
b_{ik}^{-1}
\,a_{ijk}.
\]
Using the chosen expression for \(a_{ijk}\), we get
\[
a'_{ijk}=1.
\]
Therefore
\[
\widehat g'_{ij}\,
\widehat\sigma^{\varepsilon_{ij}}(\widehat g'_{jk})
=
\widehat g'_{ik}.
\]
Thus
\[
\widehat h'_{ij}:=(\widehat g'_{ij},\varepsilon_{ij})
\]
defines a genuine \(\widehat H_{\widehat\sigma}\)-valued Čech cocycle. This cocycle defines a principal \(\widehat H_{\widehat\sigma}\)-bundle lifting \(P\).
\end{proof}

\subsection{The lifting gerbe}

Let
\[
1\longrightarrow A
\longrightarrow
\widehat H_{\widehat\sigma}
\stackrel{\pi}{\longrightarrow}
H_\sigma
\longrightarrow 1
\]
be an extension of diffeological groups, where \(A\) is abelian and central in
\(\widehat H_{\widehat\sigma}\).

Let
\[
P\to X
\]
be a principal \(H_\sigma\)-bundle.

\begin{definition}
The \emph{lifting gerbe}
\[
\mathcal G(P,\widehat H_{\widehat\sigma})
\]
is the stack over \(X\) defined as follows.

For every open subset \(U\subset X\), the groupoid
\[
\mathcal G(P,\widehat H_{\widehat\sigma})(U)
\]
has:
\begin{itemize}
\item as objects, principal
\(\widehat H_{\widehat\sigma}\)-bundles
\[
\widehat P_U\to U
\]
equipped with an isomorphism
\[
\widehat P_U/A \simeq P|_U;
\]

\item as morphisms, isomorphisms of
\(\widehat H_{\widehat\sigma}\)-bundles compatible with the projection to
\(P|_U\).
\end{itemize}
\end{definition}

\begin{remark}
The lifting gerbe is empty over \(U\) precisely when the restriction
\[
P|_U
\]
does not admit a lift to a principal
\(\widehat H_{\widehat\sigma}\)-bundle.
\end{remark}

\subsection{Band of the gerbe}

Let
\[
\alpha(P)\in H^1(X;\mathbb Z_2)
\]
denote the twisting class determined by the \(\mathbb Z_2\)-component of the
transition functions of \(P\).

The action of \(\mathbb Z_2\) on
\(\widehat H_{\widehat\sigma}\) induces an action on the central subgroup
\(A\). Hence the cocycle defining \(\alpha(P)\) determines a local system,
denoted
\[
A_\alpha.
\]

\begin{proposition}
The gerbe
\[
\mathcal G(P,\widehat H_{\widehat\sigma})
\]
is banded by the local system \(A_\alpha\).
\end{proposition}

\begin{proof}
Let
\[
\widehat P_U
\]
be a local lift over an open subset \(U\subset X\).

An automorphism of \(\widehat P_U\) inducing the identity on
\[
P|_U
\]
is given by fiberwise multiplication by an element of the kernel of
\[
\pi:
\widehat H_{\widehat\sigma}\to H_\sigma.
\]
Since the extension is central, this kernel is canonically identified with
\(A\). Therefore:
\[
\operatorname{Aut}_{P|_U}(\widehat P_U)
\simeq
A(U).
\]

Now consider two local trivializations over \(U_i\) and \(U_j\). Their
transition functions contain a \(\mathbb Z_2\)-component
\[
\varepsilon_{ij}.
\]
The action of \(\varepsilon_{ij}\) on
\(\widehat H_{\widehat\sigma}\) restricts to an automorphism of the subgroup
\(A\). Hence the identifications of the local automorphism sheaves are twisted
by the cocycle
\[
(\varepsilon_{ij}).
\]

Consequently, the automorphism sheaves glue into the local system
\[
A_\alpha.
\]
Thus the gerbe is banded by \(A_\alpha\).
\end{proof}

\subsection{Čech description of the obstruction class}

Choose a good open cover
\[
(U_i)_i
\]
of \(X\), together with transition functions
\[
(g_{ij},\varepsilon_{ij})
:
U_{ij}\to H_\sigma.
\]

Choose local lifts
\[
\widehat g_{ij}:U_{ij}\to \widehat H_{\widehat\sigma}
\]
of these transition functions.

\begin{proposition}
The functions
\[
a_{ijk}
=
\widehat g_{ij}
\widehat\sigma^{\varepsilon_{ij}}(\widehat g_{jk})
\widehat g_{ik}^{-1}
\]
take values in \(A\) and define a twisted Čech \(2\)-cocycle with coefficients
in the local system \(A_\alpha\).
\end{proposition}

\begin{proof}
Applying the projection
\[
\pi:\widehat H_{\widehat\sigma}\to H_\sigma
\]
to \(a_{ijk}\) gives
\[
(g_{ij},\varepsilon_{ij})
(g_{jk},\varepsilon_{jk})
(g_{ik},\varepsilon_{ik})^{-1}.
\]
Since the transition functions satisfy the cocycle condition in
\(H_\sigma\), this product is equal to the identity. Therefore
\[
a_{ijk}\in \ker(\pi)=A.
\]

The cocycle identity on quadruple intersections follows from associativity in
\(\widehat H_{\widehat\sigma}\). More precisely, comparing the two possible
ways of multiplying four lifted transition functions gives:
\[
a_{jkl}\,
a_{ijl}^{-1}\,
a_{ijk}\,
{}^{\varepsilon_{ij}}a_{ikl}^{-1}
=
1,
\]
which is exactly the twisted Čech cocycle condition for coefficients in the
local system \(A_\alpha\).
\end{proof}

\begin{theorem}
The cohomology class
\[
[a]\in H^2(X;A_\alpha)
\]
represented by the cocycle \((a_{ijk})\) is the class of the lifting gerbe
\[
\mathcal G(P,\widehat H_{\widehat\sigma}).
\]
\end{theorem}

\begin{proof}
The gerbe is locally non-empty because local lifts have been chosen over each
\(U_i\).

Over double intersections \(U_{ij}\), the local lifts are compared by the
transition functions
\[
\widehat g_{ij}.
\]

Over triple intersections \(U_{ijk}\), the failure of these comparison maps to
satisfy the cocycle condition is measured exactly by
\[
a_{ijk}.
\]

Changing the local lifts modifies \((a_{ijk})\) by a twisted Čech
coboundary. Hence the associated cohomology class depends only on the gerbe.

This is precisely the standard Čech construction of the class of an
\(A_\alpha\)-banded gerbe.
\end{proof}

\subsection{Twisted Dixmier--Douady class}

Assume now
\[
A=U(1).
\]

The local system \(U(1)_\alpha\) fits into the twisted exponential sequence
\[
0
\longrightarrow
\mathbb Z_\alpha
\longrightarrow
\mathbb R_\alpha
\longrightarrow
U(1)_\alpha
\longrightarrow 0.
\]

Since \(\mathbb R_\alpha\) is fine, the associated connecting morphism yields
an isomorphism
\[
H^2(X;U(1)_\alpha)
\simeq
H^3(X;\mathbb Z_\alpha).
\]

\begin{definition}
The image of
\[
[a]\in H^2(X;U(1)_\alpha)
\]
under this connecting morphism is called the \emph{twisted
Dixmier--Douady class}:
\[
\operatorname{DD}_\alpha(P)
\in
H^3(X;\mathbb Z_\alpha).
\]
\end{definition}

\begin{proposition}
The following are equivalent:
\begin{enumerate}
\item
\[
\operatorname{DD}_\alpha(P)=0;
\]

\item
\[
[a]=0
\quad\text{in }H^2(X;U(1)_\alpha);
\]

\item the lifting gerbe is trivial;

\item the principal \(H_\sigma\)-bundle \(P\) admits a lift to a principal
\(\widehat H_{\widehat\sigma}\)-bundle.
\end{enumerate}
\end{proposition}

\begin{proof}
The equivalence between (1) and (2) follows from the twisted exponential
sequence.

The equivalence between (2) and (3) is the standard classification theorem for
\(U(1)_\alpha\)-banded gerbes.

Finally, the gerbe is trivial precisely when its local objects glue to a
global lifting bundle, which is statement (4).
\end{proof}

\section{Twisted \(K\)-theory associated with the spherical construction}
\label{sec:twisted-k-theory}

\subsection{Twisting data arising from the spherical construction}

Let
\[
P\to X
\]
be a principal \(H_\sigma\)-bundle, where
\[
H_\sigma = G\rtimes_\sigma \mathbb Z_2.
\]

Associated with \(P\), the preceding sections constructed:

\begin{itemize}
\item a class
\[
\alpha(P)\in H^1(X;\mathbb Z_2),
\]
describing the induced \(\mathbb Z_2\)-local system;

\item a secondary obstruction class
\[
\beta(P)\in H^2(X;A_{\alpha(P)}),
\]
associated with the lifting gerbe.
\end{itemize}

Assume now
\[
A=U(1).
\]

Then the twisted exponential sequence
\[
0
\longrightarrow
\mathbb Z_{\alpha(P)}
\longrightarrow
\mathbb R_{\alpha(P)}
\longrightarrow
U(1)_{\alpha(P)}
\longrightarrow 0
\]
induces an isomorphism
\[
H^2(X;U(1)_{\alpha(P)})
\simeq
H^3(X;\mathbb Z_{\alpha(P)}).
\]

Hence the gerbe class determines a twisted Dixmier--Douady class
\[
\operatorname{DD}_{\alpha(P)}(P)
\in
H^3(X;\mathbb Z_{\alpha(P)}).
\]

\begin{definition}
The pair
\[
\tau(P)
:=
(\alpha(P),\operatorname{DD}_{\alpha(P)}(P))
\]
is called the \emph{spherical twisting datum} associated with the bundle
\[
P\to X.
\]
\end{definition}

\begin{remark}
The class \(\alpha(P)\) determines a local coefficient system, while
\(\operatorname{DD}_{\alpha(P)}(P)\) measures the obstruction to lifting the
bundle through the central extension
\[
1\to U(1)\to \widehat H_{\widehat\sigma}\to H_\sigma\to 1.
\]
\end{remark}

\subsection{Twisted \(K\)-theory}

Let \(X\) be a diffeological space, endowed with its \(D\)-topology.

We denote by
\[
K^\bullet(X)
\]
the complex topological \(K\)-theory of the associated topological space.

\begin{definition}
The \emph{\(\tau(P)\)-twisted \(K\)-theory groups} are the twisted
\(K\)-groups associated with the twist
\[
\tau(P)
=
(\alpha(P),\operatorname{DD}_{\alpha(P)}(P)).
\]

They are denoted by
\[
K^\bullet_{\tau(P)}(X)
\quad\text{or equivalently}\quad
K^\bullet_{\alpha(P),\beta(P)}(X).
\]
\end{definition}

\begin{remark}
One may formulate these twisted groups using any of the standard equivalent
models of twisted \(K\)-theory:
\begin{itemize}
\item bundles of Fredholm operators,
\item bundles of \(C^*\)-algebras,
\item modules over a bundle gerbe,
\item or parametrized \(K\)-theory spectra.
\end{itemize}

In the present work, only the existence and functoriality of the twisted theory
are needed.
\end{remark}

\begin{remark}
If
\[
\alpha(P)=0,
\]
then the local coefficient system becomes trivial and one recovers ordinary
gerbe-twisted \(K\)-theory classified by a degree-three integral class.
\end{remark}

\begin{remark}
If
\[
\beta(P)=0,
\]
then the twisting reduces to the \(\mathbb Z_2\)-local system determined by
\(\alpha(P)\).
\end{remark}

\subsection{Functoriality}

Let
\[
f:Y\longrightarrow X
\]
be a smooth map of diffeological spaces.

\begin{proposition}
The spherical twisting datum is functorial:
\[
f^*\tau(P)
=
(f^*\alpha(P),f^*\operatorname{DD}_{\alpha(P)}(P)).
\]

Consequently, there is a natural pull-back morphism
\[
f^*:
K^\bullet_{\tau(P)}(X)
\longrightarrow
K^\bullet_{f^*\tau(P)}(Y).
\]
\end{proposition}

\begin{proof}
The class \(\alpha(P)\) is represented by a Čech \(1\)-cocycle with values in
\(\mathbb Z_2\), while the gerbe class is represented by a twisted Čech
\(2\)-cocycle with coefficients in \(U(1)_{\alpha(P)}\).

Pulling back the cocycles along \(f\) produces representatives of the pull-back
classes. Functoriality of twisted \(K\)-theory with respect to twists then
yields the induced morphism.
\end{proof}

\subsection{Universal twisting classes}

Consider the universal principal bundle
\[
E^{\mathrm{sph}}(G)
\longrightarrow
B^{\mathrm{sph}}_\sigma(G).
\]

The associated universal twisting classes are
\[
\alpha_{\mathrm{univ}}
\in
H^1(B^{\mathrm{sph}}_\sigma(G);\mathbb Z_2)
\]
and
\[
\beta_{\mathrm{univ}}
\in
H^2(B^{\mathrm{sph}}_\sigma(G);
U(1)_{\alpha_{\mathrm{univ}}}).
\]

Their image under the twisted exponential sequence defines
\[
\operatorname{DD}_{\mathrm{univ}}
\in
H^3(B^{\mathrm{sph}}_\sigma(G);
\mathbb Z_{\alpha_{\mathrm{univ}}}).
\]

\begin{definition}
The pair
\[
\tau_{\mathrm{univ}}
=
(\alpha_{\mathrm{univ}},\operatorname{DD}_{\mathrm{univ}})
\]
is called the \emph{universal spherical \(K\)-theory twist}.
\end{definition}

\begin{theorem}
Let
\[
P\to X
\]
be a principal \(H_\sigma\)-bundle classified by
\[
f_P:X\to B^{\mathrm{sph}}_\sigma(G).
\]

Then
\[
\tau(P)
=
f_P^*(\tau_{\mathrm{univ}}).
\]
\end{theorem}

\begin{proof}
The equality for the class \(\alpha(P)\) follows from the functoriality of the
principal \(\mathbb Z_2\)-bundle defining the local system.

The equality for the gerbe class follows from functoriality of the lifting
gerbe construction and of the associated twisted Čech cocycles.

Applying the twisted exponential sequence yields the corresponding identity for
the twisted Dixmier--Douady classes.
\end{proof}

\subsection{Geometric representatives}

We briefly describe geometric representatives of twisted \(K\)-classes.

Let
\[
(U_i)_i
\]
be an open cover of \(X\), and let
\[
(a_{ijk})
\]
be a twisted gerbe cocycle representing \(\beta(P)\).

\begin{definition}
A \emph{\(\tau(P)\)-twisted vector bundle} consists of:
\begin{itemize}
\item local vector bundles
\[
E_i\to U_i,
\]

\item local isomorphisms
\[
\varphi_{ij}:
E_j|_{U_{ij}}
\longrightarrow
E_i|_{U_{ij}},
\]
\end{itemize}
such that on triple intersections
\[
\varphi_{ij}\circ\varphi_{jk}
=
a_{ijk}\,\varphi_{ik},
\]
where the coefficient \(a_{ijk}\) acts through the
\(U(1)_{\alpha(P)}\)-band of the lifting gerbe.
\end{definition}

\begin{remark}
These objects are more properly viewed as modules over the lifting gerbe
associated with the bundle \(P\).
\end{remark}

\begin{proposition}
Stable isomorphism classes of \(\tau(P)\)-twisted vector bundles determine
elements of
\[
K^0_{\tau(P)}(X).
\]
\end{proposition}

\begin{proof}
The twisted cocycle relation shows that the local bundles do not glue to an
ordinary vector bundle, but rather to a module over the gerbe determined by
\(\beta(P)\).

Direct sum defines a commutative monoid, and Grothendieck completion yields the
degree-zero twisted \(K\)-group.
\end{proof}

\subsection{Untwisted case}

\begin{theorem}
If
\[
\alpha(P)=0
\qquad\text{and}\qquad
\beta(P)=0,
\]
then there is a canonical isomorphism
\[
K^\bullet_{\tau(P)}(X)
\simeq
K^\bullet(X).
\]
\end{theorem}

\begin{proof}
If \(\alpha(P)=0\), the local coefficient system is trivial.

If \(\beta(P)=0\), the lifting gerbe is trivializable. Hence twisted vector
bundles become ordinary vector bundles after choosing a trivialization of the
gerbe.

Therefore the twisted theory reduces to ordinary complex topological
\(K\)-theory.
\end{proof}

\subsection{Summary}

The spherical Milnor construction naturally produces:
\[
\alpha(P)\in H^1(X;\mathbb Z_2),
\qquad
\beta(P)\in H^2(X;U(1)_{\alpha(P)}),
\]
and therefore a twisted Dixmier--Douady class
\[
\operatorname{DD}_{\alpha(P)}(P)
\in
H^3(X;\mathbb Z_{\alpha(P)}).
\]

These classes define canonical twists of topological \(K\)-theory associated
with spherical classifying spaces and their lifting gerbes.


\section{Descent of \(G\)-invariant geometric structures}
\label{sec:descent}

\subsection{General principle}

Let
\[
\pi:
E^{\mathrm{sph}}(G)
\longrightarrow
B^{\mathrm{sph}}(G)
=
E^{\mathrm{sph}}(G)/G
\]
be the quotient projection.

Since the action of \(G\) on \(E^{\mathrm{sph}}(G)\) is smooth, the quotient
carries the quotient diffeology associated with the subduction \(\pi\).

A geometric object defined on
\[
E^{\mathrm{sph}}(G)
\]
descends to the quotient only if it is compatible with the fibers of \(\pi\).

\subsection{Vertical tangent directions}

Let
\[
T E^{\mathrm{sph}}(G)
\]
denote the chosen tangent pseudo-bundle.

\begin{definition}
The \emph{vertical pseudo-bundle} associated with the quotient map \(\pi\) is
\[
\mathcal V
:=
\ker(d\pi)
\subset
T E^{\mathrm{sph}}(G).
\]
\end{definition}

\begin{remark}
The fibers of \(\mathcal V\) consist of tangent directions generated by the
\(G\)-orbits of the action on
\[
E^{\mathrm{sph}}(G).
\]
\end{remark}

\begin{remark}
When \(G\) is a finite-dimensional Lie group, the vertical directions are
generated by the usual fundamental vector fields. In the present diffeological
framework, especially for infinite-dimensional groups, it is more robust to use
the intrinsic definition
\[
\mathcal V=\ker(d\pi).
\]
\end{remark}

\subsection{Basic differential forms}

Let
\[
\Omega^k(E^{\mathrm{sph}}(G))
\]
denote the space of diffeological differential \(k\)-forms.

\begin{definition}
A differential form
\[
\omega\in \Omega^k(E^{\mathrm{sph}}(G))
\]
is called \emph{\(G\)-basic} if the following two conditions hold:

\begin{enumerate}
\item (\emph{\(G\)-invariance})
for every \(g\in G\),
\[
R_g^*\omega=\omega,
\]
where
\[
R_g:E^{\mathrm{sph}}(G)\to E^{\mathrm{sph}}(G)
\]
denotes the action of \(g\);

\item (\emph{horizontality})
for every point \(p\in E^{\mathrm{sph}}(G)\),
\[
\omega_p(v_1,\dots,v_k)=0
\]
whenever at least one vector \(v_i\) belongs to
\[
\mathcal V_p.
\]
\end{enumerate}
\end{definition}

\begin{remark}
The horizontality condition means precisely that \(\omega\) vanishes along the
fibers of the quotient projection.
\end{remark}

\subsection{Descent theorem}

\begin{proposition}
Let
\[
\omega\in \Omega^k(E^{\mathrm{sph}}(G))
\]
be a \(G\)-basic differential form.

Then there exists a unique differential form
\[
\bar\omega
\in
\Omega^k(B^{\mathrm{sph}}(G))
\]
such that
\[
\pi^*\bar\omega=\omega.
\]
\end{proposition}

\begin{proof}
Since
\[
\pi:
E^{\mathrm{sph}}(G)\to B^{\mathrm{sph}}(G)
\]
is a subduction, a differential form on the quotient is uniquely determined by
its pull-back to the total space.

We must therefore prove that the family of local forms defining \(\omega\) is
constant along the equivalence relation generated by the \(G\)-action.

Let
\[
p,q\in E^{\mathrm{sph}}(G)
\]
such that
\[
\pi(p)=\pi(q).
\]
Then there exists
\[
g\in G
\]
such that
\[
q=R_g(p).
\]

By \(G\)-invariance,
\[
R_g^*\omega=\omega.
\]
Hence the values of \(\omega\) at \(p\) and \(q\) correspond under the action.

Moreover, if one tangent vector belongs to
\[
\ker(d\pi),
\]
the horizontality condition implies that the value of \(\omega\) vanishes.
Therefore \(\omega\) depends only on tangent directions projected to the
quotient.

Consequently, \(\omega\) defines uniquely a differential form
\[
\bar\omega
\]
on
\[
B^{\mathrm{sph}}(G)
\]
whose pull-back is \(\omega\).

Uniqueness follows from the fact that \(\pi\) is a subduction.
\end{proof}

\subsection{Basic de Rham complex}

\begin{definition}
We denote by
\[
\Omega^\bullet_{\mathrm{bas}}(E^{\mathrm{sph}}(G))
\]
the graded vector space of \(G\)-basic differential forms.
\end{definition}

\begin{proposition}
The exterior differential preserves basic forms:
\[
d:
\Omega^k_{\mathrm{bas}}(E^{\mathrm{sph}}(G))
\longrightarrow
\Omega^{k+1}_{\mathrm{bas}}(E^{\mathrm{sph}}(G)).
\]
\end{proposition}

\begin{proof}
The exterior differential commutes with pull-backs, hence preserves
\(G\)-invariance.

Moreover, if one argument is vertical, Cartan's formula for the exterior
differential shows that \(d\omega\) still vanishes whenever one entry belongs to
\(\mathcal V\). Therefore horizontality is preserved.
\end{proof}

\begin{proposition}
There is a canonical identification
\[
\Omega^\bullet(B^{\mathrm{sph}}(G))
\simeq
\Omega^\bullet_{\mathrm{bas}}(E^{\mathrm{sph}}(G)).
\]
\end{proposition}

\subsection{Descent of the barycentric metric}

Recall that the barycentric metric on \(E^{\mathrm{sph}}(G)\) is locally given by
\[
g
=
\sum_i dx_i^2
+
\sum_i x_i^2 \langle \cdot,\cdot\rangle_G^{(i)}.
\]

\begin{proposition}
Assume that the metric \(\langle\cdot,\cdot\rangle_G\) on \(G\) is left-invariant. Then the barycentric metric \(g\) is \(G\)-invariant.
\end{proposition}

\begin{proof}
The \(G\)-action on \(E^{\mathrm{sph}}(G)\) is
\[
h\cdot (x_i,g_i)=(x_i,hg_i).
\]
The coordinates \(x_i\) are unchanged by this action, hence the term
\[
\sum_i dx_i^2
\]
is invariant.

For the group part, left multiplication by \(h\) sends \(g_i\) to \(hg_i\). Since the metric on \(G\) is left-invariant,
\[
\langle d(hg_i)v,d(hg_i)w\rangle_G
=
\langle v,w\rangle_G.
\]
Thus each term
\[
x_i^2\langle\cdot,\cdot\rangle_G^{(i)}
\]
is invariant. Hence \(g\) is \(G\)-invariant.
\end{proof}

\begin{remark}
The metric is not automatically basic as a tensor on the full tangent pseudo-bundle, because it may have nonzero components along the \(G\)-orbits. What descends canonically is the induced metric on the quotient horizontal directions, once the vertical directions are quotiented out.
\end{remark}

\begin{definition}
Let
\[
\mathcal V:=\ker(d\pi)
\]
be the vertical pseudo-bundle of the quotient map. The quotient tangent pseudo-bundle is
\[
TB^{\mathrm{sph}}(G)
\simeq
TE^{\mathrm{sph}}(G)/\mathcal V.
\]
\end{definition}

\begin{proposition}
If the barycentric metric is \(G\)-invariant and its vertical kernel is compatible with the quotient, it induces a metric-type structure
\[
\bar g
\]
on \(B^{\mathrm{sph}}(G)\).
\end{proposition}

\begin{proof}
Since \(g\) is \(G\)-invariant, its value is constant along \(G\)-orbits. To define a tensor on the quotient, one must verify that changing representatives by vertical vectors does not change the value. This is equivalent to the vertical directions being killed or consistently quotiented in the induced horizontal pseudo-bundle.

Under the stated compatibility condition, the formula
\[
\bar g(d\pi(v),d\pi(w)):=g(v,w)
\]
is independent of the choice of lifts \(v,w\). Smoothness follows from the quotient diffeology.
\end{proof}

\subsection{Descent of differential forms}

\begin{proposition}
Let
\[
\omega\in\Omega^k_{\mathrm{bas}}(E^{\mathrm{sph}}(G))
\]
be a \(G\)-basic differential form.

Then
\[
d\omega
\]
is again \(G\)-basic.

Moreover, if
\[
\bar\omega\in\Omega^k(B^{\mathrm{sph}}(G))
\]
denotes the descended form satisfying
\[
\pi^*\bar\omega=\omega,
\]
then
\[
d\bar\omega
=
\overline{d\omega}.
\]
\end{proposition}

\begin{proof}
We first prove \(G\)-invariance.

Let
\[
g\in G.
\]
Since pull-back commutes with the exterior differential,
\[
R_g^*(d\omega)
=
d(R_g^*\omega).
\]
Because \(\omega\) is \(G\)-invariant,
\[
R_g^*\omega=\omega,
\]
hence
\[
R_g^*(d\omega)=d\omega.
\]

Thus \(d\omega\) is \(G\)-invariant.

We now prove horizontality.

Let
\[
v_0,\dots,v_k
\in
T_pE^{\mathrm{sph}}(G)
\]
with one vector, say \(v_0\), belonging to the vertical subspace
\[
\mathcal V_p=\ker(d\pi_p).
\]

Since \(\omega\) is horizontal, it vanishes whenever one argument is vertical.
Using the local expression of the exterior differential on plots, every term
appearing in
\[
(d\omega)_p(v_0,\dots,v_k)
\]
either contains evaluation of \(\omega\) on a family including a vertical
vector, or derivatives along directions tangent to \(G\)-orbits of functions
already vanishing identically by horizontality and invariance.

Therefore
\[
(d\omega)_p(v_0,\dots,v_k)=0.
\]

Hence \(d\omega\) is horizontal.

Consequently,
\[
d\omega
\]
is \(G\)-basic.

Finally,
\[
\pi^*(d\bar\omega)
=
d(\pi^*\bar\omega)
=
d\omega.
\]
Since \(d\omega\) is basic, it descends uniquely to a form on
\[
B^{\mathrm{sph}}(G).
\]
By uniqueness of descent,
\[
d\bar\omega
=
\overline{d\omega}.
\]
\end{proof}

\subsection{Descent of characteristic forms}

We now discuss characteristic forms associated with principal connections.

\begin{remark}
In the fully general diffeological setting, the existence of a Lie algebra,
connection \(1\)-forms, and curvature forms requires additional assumptions on
the diffeological group \(G\).

Accordingly, in this subsection we assume that \(G\) is a regular
diffeological Lie group admitting:
\begin{itemize}
\item a well-defined Lie algebra \(\mathfrak g\),
\item a smooth adjoint action,
\item and a theory of principal connections with curvature.
\end{itemize}
Typical examples include finite-dimensional Lie groups, Fréchet--Lie groups,
loop groups, and current groups.
\end{remark}

Let
\[
\pi:
E^{\mathrm{sph}}(G)\to B^{\mathrm{sph}}(G)
\]
be the universal principal \(G\)-bundle.

Let
\[
\Theta
\in
\Omega^1(E^{\mathrm{sph}}(G),\mathfrak g)
\]
be a principal connection form.

\begin{remark}
The connection form itself does not descend to the quotient, since it is not
horizontal with respect to the principal action.
\end{remark}

Its curvature is defined by
\[
F_\Theta
=
d\Theta
+
\frac12[\Theta,\Theta].
\]

\begin{proposition}
The curvature form
\[
F_\Theta
\]
is horizontal and \(G\)-equivariant:
\[
R_g^*F_\Theta
=
\operatorname{Ad}_{g^{-1}}F_\Theta.
\]
\end{proposition}

\begin{proof}
This is the standard structural identity for principal connections and follows
from the equivariance of \(\Theta\) together with the compatibility of the
exterior differential and the Lie bracket with pull-backs.
\end{proof}

Let
\[
P:\mathfrak g^{\otimes r}\to\mathbb R
\]
be an Ad-invariant symmetric polynomial.

\begin{proposition}
The characteristic form
\[
P(F_\Theta,\dots,F_\Theta)
\]
is \(G\)-basic.

Consequently, it descends uniquely to a differential form on
\[
B^{\mathrm{sph}}(G).
\]
\end{proposition}

\begin{proof}
Since \(F_\Theta\) is horizontal, the multilinear form
\[
P(F_\Theta,\dots,F_\Theta)
\]
vanishes whenever one argument is vertical. Hence it is horizontal.

Moreover,
\[
R_g^*F_\Theta
=
\operatorname{Ad}_{g^{-1}}F_\Theta.
\]
Therefore,
\[
R_g^*P(F_\Theta,\dots,F_\Theta)
=
P(
\operatorname{Ad}_{g^{-1}}F_\Theta,
\dots,
\operatorname{Ad}_{g^{-1}}F_\Theta
).
\]

By Ad-invariance of \(P\),
\[
P(
\operatorname{Ad}_{g^{-1}}X_1,
\dots,
\operatorname{Ad}_{g^{-1}}X_r
)
=
P(X_1,\dots,X_r),
\]
hence
\[
R_g^*P(F_\Theta,\dots,F_\Theta)
=
P(F_\Theta,\dots,F_\Theta).
\]

Thus the characteristic form is \(G\)-invariant and horizontal, hence
\(G\)-basic.

The descent theorem for basic forms therefore produces a unique form on
\[
B^{\mathrm{sph}}(G)
\]
whose pull-back is
\[
P(F_\Theta,\dots,F_\Theta).
\]
\end{proof}
\subsection{Descent of Hodge-type operators}

In this subsection, we assume that we are working in a framework where the
Hodge-type structures constructed on
\[
E^{\mathrm{sph}}(G)
\]
are well-defined. More precisely, we assume that:
\begin{itemize}
\item the horizontal differential complex is defined;
\item a metric pairing on horizontal forms is available;
\item a formal adjoint
\[
\delta
\]
of the horizontal exterior differential exists;
\item the quotient
\[
B^{\mathrm{sph}}(G)
\]
carries the descended horizontal differential calculus.
\end{itemize}

Under these assumptions, the descent question is purely equivariant.

\begin{definition}
The basic horizontal complex is
\[
\Omega^\bullet_{\mathrm{bas}}(E^{\mathrm{sph}}(G)).
\]
By the descent theorem for basic forms, it identifies with the de Rham complex
of the quotient:
\[
\Omega^\bullet_{\mathrm{bas}}(E^{\mathrm{sph}}(G))
\simeq
\Omega^\bullet(B^{\mathrm{sph}}(G)).
\]
\end{definition}

\begin{definition}
We say that the Hodge-type structure is \(G\)-stable if
\[
\delta\bigl(
\Omega^{k+1}_{\mathrm{bas}}(E^{\mathrm{sph}}(G))
\bigr)
\subset
\Omega^k_{\mathrm{bas}}(E^{\mathrm{sph}}(G))
\]
for every \(k\).
\end{definition}

\begin{proposition}
Assume that the Hodge-type structure is \(G\)-stable. Then the formal adjoint
\[
\delta
\]
descends to an operator
\[
\bar\delta:
\Omega^{k+1}(B^{\mathrm{sph}}(G))
\longrightarrow
\Omega^k(B^{\mathrm{sph}}(G)).
\]

Consequently, the operator
\[
\bar\Delta
=
d\bar\delta+\bar\delta d
\]
defines a Hodge-type Laplacian on the quotient.
\end{proposition}

\begin{proof}
Let
\[
\bar\omega\in\Omega^{k+1}(B^{\mathrm{sph}}(G)).
\]
By the basic-form identification, its pull-back
\[
\omega:=\pi^*\bar\omega
\]
is a basic \((k+1)\)-form on
\[
E^{\mathrm{sph}}(G).
\]

Since the Hodge-type structure is \(G\)-stable,
\[
\delta\omega
\]
is again basic. Hence there exists a unique form
\[
\bar\eta\in\Omega^k(B^{\mathrm{sph}}(G))
\]
such that
\[
\pi^*\bar\eta=\delta\omega.
\]

Define
\[
\bar\delta\bar\omega:=\bar\eta.
\]

This definition is independent of choices because the pull-back map
\[
\pi^*:\Omega^\bullet(B^{\mathrm{sph}}(G))
\to
\Omega^\bullet_{\mathrm{bas}}(E^{\mathrm{sph}}(G))
\]
is injective for the quotient subduction.

Thus
\[
\bar\delta
\]
is well-defined.

Since the exterior differential preserves basic forms and commutes with
descent, the expression
\[
\bar\Delta=d\bar\delta+\bar\delta d
\]
is a well-defined operator on the quotient complex.
\end{proof}

\begin{remark}
The proposition does not assert that a formal adjoint exists in complete
generality. It only states that, once such an operator has been constructed on
the total space, its descent to the quotient is equivalent to preservation of
basic forms.
\end{remark}

\subsection{Descent of Clifford and Dirac structures}

We now assume that we are in a framework where the Clifford and Dirac
structures on
\[
E^{\mathrm{sph}}(G)
\]
are well-defined. Thus we assume that:
\begin{itemize}
\item an effective horizontal pseudo-bundle
\[
\mathcal H_E
\]
has been constructed;
\item \(\mathcal H_E\) carries a metric;
\item a Clifford algebra pseudo-bundle
\[
\mathrm{Cl}(\mathcal H_E)
\]
is defined;
\item a Clifford module
\[
\mathcal E_E
\]
has been chosen;
\item a Clifford-compatible connection
\[
\nabla^{\mathcal E}
\]
exists;
\item the associated Dirac-type operator
\[
D_E
\]
is well-defined.
\end{itemize}

\subsubsection{Equivariant Clifford data}

\begin{definition}
The Clifford data are called \(G\)-equivariant if:
\begin{enumerate}
\item the \(G\)-action lifts to \(\mathcal H_E\);
\item the metric on \(\mathcal H_E\) is \(G\)-invariant;
\item the Clifford action satisfies
\[
g\cdot(c(v)s)
=
c(g\cdot v)(g\cdot s)
\]
for all admissible \(g\in G\), \(v\in\mathcal H_E\), and
\(s\in\mathcal E_E\);
\item the connection preserves the equivariant structure.
\end{enumerate}
\end{definition}

\begin{proposition}
If the Clifford data are \(G\)-equivariant, then the horizontal pseudo-bundle,
the Clifford algebra, and the Clifford module descend to the quotient
\[
B^{\mathrm{sph}}(G).
\]
\end{proposition}

\begin{proof}
Since the \(G\)-action lifts to
\[
\mathcal H_E,
\]
one can form the quotient pseudo-bundle
\[
\mathcal H_B:=\mathcal H_E/G
\]
over
\[
B^{\mathrm{sph}}(G).
\]

The \(G\)-invariance of the metric implies that the metric descends to
\[
\mathcal H_B.
\]

The Clifford relations are preserved by the \(G\)-action because the Clifford
action is equivariant. Therefore the quotient of
\[
\mathrm{Cl}(\mathcal H_E)
\]
is naturally identified with
\[
\mathrm{Cl}(\mathcal H_B).
\]

Similarly, the equivariance of the Clifford module implies that the quotient
\[
\mathcal E_B:=\mathcal E_E/G
\]
is a Clifford module over
\[
\mathrm{Cl}(\mathcal H_B).
\]
\end{proof}

\subsubsection{Descent of the Dirac operator}

\begin{definition}
We say that the Dirac operator
\[
D_E
\]
is \(G\)-stable if it maps \(G\)-basic sections of
\[
\mathcal E_E
\]
to \(G\)-basic sections.
\end{definition}

\begin{proposition}
Assume that:
\begin{enumerate}
\item the Clifford data are \(G\)-equivariant;
\item the connection is compatible with the \(G\)-action;
\item the Dirac operator \(D_E\) is \(G\)-stable.
\end{enumerate}

Then \(D_E\) descends to a Dirac-type operator
\[
\bar D:
\Gamma(\mathcal E_B)
\longrightarrow
\Gamma(\mathcal E_B)
\]
on
\[
B^{\mathrm{sph}}(G).
\]
\end{proposition}

\begin{proof}
Let
\[
\bar s\in\Gamma(\mathcal E_B)
\]
be a section of the descended Clifford module.

By descent of the module, \(\bar s\) corresponds uniquely to a \(G\)-basic
section
\[
s=\pi^*\bar s
\]
of
\[
\mathcal E_E.
\]

Since \(D_E\) is \(G\)-stable,
\[
D_Es
\]
is again a \(G\)-basic section. Hence it descends uniquely to a section
\[
\bar t\in\Gamma(\mathcal E_B)
\]
such that
\[
\pi^*\bar t=D_Es.
\]

Define
\[
\bar D\bar s:=\bar t.
\]

This is independent of choices because the pull-back correspondence between
sections of \(\mathcal E_B\) and \(G\)-basic sections of \(\mathcal E_E\) is
injective.

Thus
\[
\bar D
\]
is a well-defined operator on the quotient.
\end{proof}

\begin{remark}
The statement is purely a descent result. It does not assert that the Dirac
operator exists automatically in arbitrary diffeological or
infinite-dimensional settings. It only asserts that, once a well-defined
Dirac-type operator exists on the total space, it descends precisely under
the stated equivariance and stability assumptions.
\end{remark}
\subsection{Pullback by classifying maps}

Let
\[
f:X\longrightarrow B^{\mathrm{sph}}(G)
\]
be a smooth classifying map. Once a geometric object has descended to \(B^{\mathrm{sph}}(G)\), it can be pulled back to \(X\).

\begin{proposition}
Every descended object on \(B^{\mathrm{sph}}(G)\), such as:
\[
\bar g,\quad
\bar\omega,\quad
\bar\Delta,\quad
\bar D,
\]
induces by pullback a corresponding object on \(X\):
\[
f^*\bar g,\quad
f^*\bar\omega,\quad
f^*\bar\Delta,\quad
f^*\bar D,
\]
whenever the corresponding pullback operation is defined in the diffeological category.
\end{proposition}

\begin{proof}
For tensors and forms, this is the ordinary pullback operation. For pseudo-bundles and operators, one pulls back the underlying pseudo-bundle and then the associated connection or operator. The construction is functorial because all objects were first defined on the quotient base \(B^{\mathrm{sph}}(G)\).
\end{proof}

\begin{remark}
This section explains the universal nature of the constructions above. The objects that are \(G\)-basic descend to the classifying base; once descended, they generate geometric structures on any diffeological space equipped with a classifying map.
\end{remark}

\section{\(G\)-stability of Hodge and Dirac structures}
\label{sec:G-stability-Hodge-Dirac}

We now analyze the behavior of Hodge-type and Dirac-type structures under the
action of \(G\). Throughout this section, we assume that we are working in a
framework where the relevant Hodge and Dirac structures are well-defined.

The guiding principle is the following: a structure on
\[
E^{\mathrm{sph}}(G)
\]
descends to
\[
B^{\mathrm{sph}}(G)=E^{\mathrm{sph}}(G)/G
\]
only if it is compatible with the quotient projection. When this compatibility
fails, the structure remains defined on the total space, and the failure of
descent is measured by a defect.

\subsection{Basic forms and the quotient complex}

Let
\[
\pi:E^{\mathrm{sph}}(G)\to B^{\mathrm{sph}}(G)
\]
be the quotient projection, and let
\[
\mathcal V:=\ker(d\pi)
\]
be the vertical pseudo-bundle.

\begin{definition}
A differential form
\[
\omega\in\Omega^k(E^{\mathrm{sph}}(G))
\]
is called \(G\)-basic if:
\begin{enumerate}
\item it is \(G\)-invariant:
\[
R_g^*\omega=\omega
\qquad
\text{for all }g\in G;
\]
\item it is horizontal:
\[
\omega(v_1,\ldots,v_k)=0
\]
whenever one of the vectors \(v_i\) belongs to \(\mathcal V\).
\end{enumerate}
\end{definition}

\begin{proposition}
There is a natural identification
\[
\Omega^\bullet(B^{\mathrm{sph}}(G))
\simeq
\Omega^\bullet_{\mathrm{bas}}(E^{\mathrm{sph}}(G)).
\]
\end{proposition}

\begin{proof}
This is the descent theorem for differential forms along the quotient
subduction
\[
\pi:E^{\mathrm{sph}}(G)\to B^{\mathrm{sph}}(G).
\]
A form on the quotient pulls back to a \(G\)-invariant horizontal form on the
total space. Conversely, a \(G\)-basic form depends only on tangent directions
modulo \(\ker(d\pi)\) and is constant along \(G\)-orbits, hence descends
uniquely to the quotient.
\end{proof}

\subsection{Stability of the Hodge complex}

Assume that the horizontal de Rham complex and the codifferential
\[
\delta:
\Omega^{k+1}(E^{\mathrm{sph}}(G))
\longrightarrow
\Omega^k(E^{\mathrm{sph}}(G))
\]
are defined.

\begin{definition}
The Hodge-type structure is called \(G\)-stable if
\[
\delta\bigl(
\Omega^{k+1}_{\mathrm{bas}}(E^{\mathrm{sph}}(G))
\bigr)
\subset
\Omega^k_{\mathrm{bas}}(E^{\mathrm{sph}}(G)).
\]
\end{definition}

\begin{proposition}
If the Hodge-type structure is \(G\)-stable, then \(\delta\) descends to an
operator
\[
\bar\delta:
\Omega^{k+1}(B^{\mathrm{sph}}(G))
\to
\Omega^k(B^{\mathrm{sph}}(G)).
\]
Consequently,
\[
\bar\Delta=d\bar\delta+\bar\delta d
\]
defines a Hodge-type Laplacian on the quotient.
\end{proposition}

\begin{proof}
Let
\[
\bar\omega\in \Omega^{k+1}(B^{\mathrm{sph}}(G)).
\]
Then
\[
\omega:=\pi^*\bar\omega
\]
is \(G\)-basic. Since the Hodge structure is \(G\)-stable,
\[
\delta\omega
\]
is again \(G\)-basic. Therefore there exists a unique form
\[
\bar\eta\in\Omega^k(B^{\mathrm{sph}}(G))
\]
such that
\[
\pi^*\bar\eta=\delta\omega.
\]
Define
\[
\bar\delta\bar\omega:=\bar\eta.
\]
This is well-defined because pull-back by a quotient subduction is injective
on descended forms. Since \(d\) preserves basic forms, the operator
\[
\bar\Delta=d\bar\delta+\bar\delta d
\]
is well-defined on the quotient.
\end{proof}

\begin{remark}
The assertion is conditional. It does not claim that a codifferential exists in
full generality; it only describes its descent once it has been constructed.
\end{remark}

\subsection{The Hodge descent defect}

If \(\delta\) does not preserve basic forms, one obtains a natural obstruction.

\begin{definition}
The \emph{Hodge descent defect} is the map
\[
\mathcal D_{\mathrm{Hodge}}
:
\Omega^{k+1}_{\mathrm{bas}}(E^{\mathrm{sph}}(G))
\longrightarrow
\Omega^k(E^{\mathrm{sph}}(G))/
\Omega^k_{\mathrm{bas}}(E^{\mathrm{sph}}(G))
\]
defined by
\[
\mathcal D_{\mathrm{Hodge}}(\omega)
=
[\delta\omega].
\]
\end{definition}

\begin{proposition}
The Hodge-type structure descends to the quotient if and only if
\[
\mathcal D_{\mathrm{Hodge}}=0.
\]
\end{proposition}

\begin{proof}
The operator \(\delta\) descends if and only if it sends every basic form to a
basic form. This condition is equivalent to saying that, for every basic form
\(\omega\), the class of \(\delta\omega\) in
\[
\Omega^k(E^{\mathrm{sph}}(G))/
\Omega^k_{\mathrm{bas}}(E^{\mathrm{sph}}(G))
\]
vanishes. This is exactly
\[
\mathcal D_{\mathrm{Hodge}}=0.
\]
\end{proof}

\subsection{\(G\)-equivariant Clifford structures}

Let
\[
\mathcal H_E\to E^{\mathrm{sph}}(G)
\]
be the effective horizontal pseudo-bundle, and assume that it carries a metric,
a Clifford algebra pseudo-bundle
\[
\mathrm{Cl}(\mathcal H_E),
\]
and a Clifford module
\[
\mathcal E_E.
\]

\begin{definition}
The Clifford data are called \(G\)-equivariant if:
\begin{enumerate}
\item the \(G\)-action lifts to \(\mathcal H_E\);
\item the metric on \(\mathcal H_E\) is \(G\)-invariant;
\item the \(G\)-action lifts to \(\mathcal E_E\);
\item Clifford multiplication satisfies
\[
g\cdot(c(v)s)
=
c(g\cdot v)(g\cdot s)
\]
for all admissible \(g\in G\), \(v\in\mathcal H_E\), and
\(s\in\mathcal E_E\).
\end{enumerate}
\end{definition}

\begin{proposition}
If the Clifford data are \(G\)-equivariant, then the horizontal pseudo-bundle,
the Clifford algebra, and the Clifford module descend to the quotient.
\end{proposition}

\begin{proof}
Since the action of \(G\) lifts to \(\mathcal H_E\), the quotient
\[
\mathcal H_B:=\mathcal H_E/G
\]
is a pseudo-bundle over
\[
B^{\mathrm{sph}}(G).
\]
The metric descends because it is \(G\)-invariant.

The Clifford relation
\[
c(v)c(w)+c(w)c(v)=-2\bar g(v,w)
\]
is preserved by the \(G\)-action. Therefore the quotient Clifford algebra is
naturally identified with
\[
\mathrm{Cl}(\mathcal H_B).
\]
Finally, the equivariant Clifford module descends to a Clifford module
\[
\mathcal E_B:=\mathcal E_E/G
\]
over the quotient.
\end{proof}

\subsection{\(G\)-stability of the Dirac operator}

Assume that a Dirac-type operator
\[
D_E:\Gamma(\mathcal E_E)\to \Gamma(\mathcal E_E)
\]
has been constructed from the Clifford data and a compatible connection.

\begin{definition}
The operator \(D_E\) is called \(G\)-equivariant if
\[
g^*(D_Es)=D_E(g^*s)
\]
for every admissible \(g\in G\) and every section \(s\).
\end{definition}

\begin{definition}
The operator \(D_E\) is called \(G\)-stable if it maps \(G\)-basic sections to
\(G\)-basic sections.
\end{definition}

\begin{proposition}
If \(D_E\) is \(G\)-stable, then it descends to an operator
\[
\bar D:\Gamma(\mathcal E_B)\to\Gamma(\mathcal E_B).
\]
\end{proposition}

\begin{proof}
A section
\[
\bar s\in\Gamma(\mathcal E_B)
\]
corresponds to a \(G\)-basic section
\[
s=\pi^*\bar s
\]
of \(\mathcal E_E\). Since \(D_E\) is \(G\)-stable,
\[
D_Es
\]
is again basic. Therefore it descends uniquely to a section
\[
\bar t\in\Gamma(\mathcal E_B).
\]
Define
\[
\bar D\bar s:=\bar t.
\]
This gives a well-defined operator on the quotient.
\end{proof}

\subsection{The Dirac equivariance defect}

If \(D_E\) is not \(G\)-equivariant, the obstruction is measured by its failure
to commute with the \(G\)-action.

\begin{definition}
For \(g\in G\), define the \emph{Dirac equivariance defect}
\[
\mathcal D_D(g)
:=
g^*D_E-D_Eg^*.
\]
Thus \(D_E\) is \(G\)-equivariant if and only if
\[
\mathcal D_D(g)=0
\]
for every \(g\in G\).
\end{definition}

\begin{remark}
If \(G\) is a regular diffeological Lie group and the action is generated by
its Lie algebra, one may also define an infinitesimal defect
\[
\mathcal D_D(\xi)
=
[\mathcal L_{\xi^\#},D_E],
\qquad
\xi\in\mathfrak g.
\]
This infinitesimal expression is used only under the additional assumption
that fundamental vector fields and their Lie derivatives are well-defined.
\end{remark}

\subsection{Pullback to a classified principal bundle}

Let
\[
f:X\to B^{\mathrm{sph}}(G)
\]
be a classifying map, and let
\[
P=f^*E^{\mathrm{sph}}(G)
\]
be the associated principal \(G\)-bundle.

Even when a geometric object on \(E^{\mathrm{sph}}(G)\) does not descend to
\(B^{\mathrm{sph}}(G)\), it pulls back to the total space \(P\).

\begin{proposition}
Let \(\mathcal O\) be a geometric object on \(E^{\mathrm{sph}}(G)\) for which
pullback is defined. Then
\[
\widetilde{\mathcal O}:=\tilde f^*\mathcal O
\]
is a well-defined object on \(P\), where
\[
\tilde f:P\to E^{\mathrm{sph}}(G)
\]
is the canonical projection to the second factor.
\end{proposition}

\begin{proof}
The pullback bundle is the fiber product
\[
P
=
\{(x,e)\in X\times E^{\mathrm{sph}}(G)\mid f(x)=\pi(e)\}.
\]
The map
\[
\tilde f:P\to E^{\mathrm{sph}}(G),
\qquad
(x,e)\mapsto e,
\]
is smooth. Therefore every pullback-compatible object on
\(E^{\mathrm{sph}}(G)\) pulls back to \(P\).
\end{proof}

\subsection{Defect on the pulled-back bundle}

Let
\[
D_P:=\tilde f^*D_E
\]
be the pulled-back Dirac-type operator on \(P\).

\begin{definition}
For \(g\in G\), the \emph{pulled-back Dirac defect} is
\[
\mathcal D_{D,P}(g)
:=
g^*D_P-D_Pg^*.
\]
\end{definition}

\begin{proposition}
The operator \(D_P\) descends from \(P\) to \(X=P/G\) if and only if it maps
basic sections to basic sections. Equivalently, its descent obstruction is the
failure of
\[
D_P
\]
to preserve the basic subspace.
\end{proposition}

\begin{proof}
Sections over \(X=P/G\) correspond to \(G\)-basic sections over \(P\). Hence an
operator on \(P\) descends to an operator on \(X\) exactly when it sends basic
sections to basic sections.

If this condition holds, define the descended operator by applying \(D_P\) to
the basic lift of a quotient section and descending the result. If the
condition fails, such a quotient operator cannot be defined.
\end{proof}

\begin{remark}
If one works in a Lie-regular framework, preservation of basic sections may be
tested infinitesimally by commutators with vertical Lie derivatives. Outside
that framework, the basic-section criterion is the safer formulation.
\end{remark}

\subsection{Summary}

There are three regimes.

\begin{enumerate}
\item If a structure is basic, it descends directly to
\[
B^{\mathrm{sph}}(G).
\]

\item If a structure is equivariant but not horizontal, it may still induce
associated quotient data, for example through curvature, associated bundles,
or basic components.

\item If a structure is not equivariant or does not preserve basic sections, it
does not descend. Nevertheless, it pulls back to the total space of any
classified principal bundle
\[
P=f^*E^{\mathrm{sph}}(G),
\]
and the failure of descent is measured by defects such as
\[
\mathcal D_{\mathrm{Hodge}},
\qquad
\mathcal D_D.
\]
\end{enumerate}

\section{Characteristic classes associated with equivariance defects}
\label{sec:defect-classes}

In this section, we associate secondary characteristic classes with the
failure of equivariance of geometric operators on
\[
E^{\mathrm{sph}}(G).
\]
The philosophy is analogous to Chern--Weil theory: instead of using the
curvature of a connection, we use the failure of a geometric operator to
commute with the \(G\)-action.

Throughout this section, we assume that:
\begin{enumerate}
\item \(G\) is a regular diffeological Lie group admitting a Lie algebra
\(\mathfrak g\);
\item the action of \(G\) on \(E^{\mathrm{sph}}(G)\) admits fundamental vector
fields;
\item Lie derivatives of the geometric objects under consideration are
well-defined.
\end{enumerate}

\subsection{Infinitesimal equivariance defect}

Let
\[
\pi:E^{\mathrm{sph}}(G)\to B^{\mathrm{sph}}(G)
\]
be the universal principal \(G\)-bundle.

Let
\[
D:\Gamma(\mathcal E)\to\Gamma(\mathcal E)
\]
be a Dirac-type operator acting on a Clifford module
\[
\mathcal E\to E^{\mathrm{sph}}(G).
\]

For each
\[
\xi\in\mathfrak g,
\]
let
\[
\xi^\#
\]
denote the corresponding fundamental vector field on
\[
E^{\mathrm{sph}}(G).
\]

\begin{definition}
The \emph{infinitesimal equivariance defect} of \(D\) is the map
\[
\mathcal D_D:\mathfrak g
\longrightarrow
\operatorname{End}(\Gamma(\mathcal E))
\]
defined by
\[
\mathcal D_D(\xi)
:=
[\mathcal L_{\xi^\#},D].
\]
\end{definition}

\begin{remark}
The defect vanishes identically if and only if \(D\) commutes with all
infinitesimal generators of the \(G\)-action.
\end{remark}

\begin{proposition}
The map
\[
\mathcal D_D
\]
is linear in \(\xi\).
\end{proposition}

\begin{proof}
The assignment
\[
\xi\mapsto \xi^\#
\]
is linear, and the Lie derivative depends linearly on the vector field.
Since the commutator is bilinear,
\[
\mathcal D_D(a\xi+b\eta)
=
a\mathcal D_D(\xi)
+
b\mathcal D_D(\eta).
\]
\end{proof}

\subsection{Curvature of the defect}

The defect behaves formally like a connection \(1\)-form on the trivial bundle
with fiber \(\operatorname{End}(\Gamma(\mathcal E))\).

\begin{definition}
The \emph{defect curvature} is the alternating bilinear map
\[
\mathcal F_D:
\mathfrak g\times\mathfrak g
\longrightarrow
\operatorname{End}(\Gamma(\mathcal E))
\]
defined by
\[
\mathcal F_D(\xi,\eta)
:=
[\mathcal D_D(\xi),\mathcal D_D(\eta)]
-
\mathcal D_D([\xi,\eta]).
\]
\end{definition}

\begin{remark}
This is formally the curvature associated with the ``connection''
\[
\mathcal D_D.
\]
It measures the failure of \(\mathcal D_D\) to define a representation of the
Lie algebra \(\mathfrak g\).
\end{remark}

\begin{proposition}
The map
\[
\mathcal F_D
\]
is bilinear and alternating.
\end{proposition}

\begin{proof}
Bilinearity follows from bilinearity of the commutator and linearity of
\(\mathcal D_D\).

For antisymmetry,
\[
\mathcal F_D(\eta,\xi)
=
[\mathcal D_D(\eta),\mathcal D_D(\xi)]
-
\mathcal D_D([\eta,\xi]).
\]
Since
\[
[A,B]=-[B,A]
\]
and
\[
[\eta,\xi]=-[\xi,\eta],
\]
we obtain
\[
\mathcal F_D(\eta,\xi)
=
-\mathcal F_D(\xi,\eta).
\]
\end{proof}

\subsection{Invariant polynomials}

Let
\[
\mathcal A
\subset
\operatorname{End}(\Gamma(\mathcal E))
\]
be an algebra containing the image of
\[
\mathcal F_D.
\]

\begin{definition}
An \emph{invariant polynomial} on \(\mathcal A\) is a symmetric multilinear map
\[
P:\mathcal A^k\to\mathbb R
\]
such that
\[
P(gAg^{-1},\ldots,gAg^{-1})
=
P(A,\ldots,A)
\]
whenever the conjugation action is defined.
\end{definition}

Typical examples are traces of products whenever a trace functional exists.

\subsection{Characteristic cochains}

\begin{definition}
Let \(P\) be an invariant polynomial of degree \(k\). Define the alternating
\(2k\)-cochain
\[
\omega_P
\in
\Lambda^{2k}\mathfrak g^*
\]
by
\[
\omega_P(\xi_1,\ldots,\xi_{2k})
=
\sum_{\sigma}
\operatorname{sgn}(\sigma)\,
P\bigl(
\mathcal F_D(\xi_{\sigma(1)},\xi_{\sigma(2)}),
\ldots,
\mathcal F_D(\xi_{\sigma(2k-1)},\xi_{\sigma(2k)})
\bigr),
\]
where the sum runs over pairings of
\[
\{1,\ldots,2k\}.
\]
\end{definition}

\begin{remark}
This construction is the direct analogue of the Chern--Weil construction,
with the defect curvature replacing the curvature of a connection.
\end{remark}

\subsection{Closedness}

\begin{proposition}
Assume that:
\begin{enumerate}
\item the defect satisfies the Bianchi-type identity
\[
d_{\mathrm{CE}}\mathcal F_D=0
\]
in Chevalley--Eilenberg cohomology;
\item \(P\) is invariant.
\end{enumerate}
Then
\[
d_{\mathrm{CE}}\omega_P=0.
\]
\end{proposition}

\begin{proof}
The proof is formally identical to the classical Chern--Weil argument.
The Chevalley--Eilenberg differential of
\[
\omega_P
\]
is obtained by differentiating the entries of
\[
P(\mathcal F_D,\ldots,\mathcal F_D).
\]
Using the Bianchi identity and invariance of \(P\), all terms cancel.
Hence
\[
d_{\mathrm{CE}}\omega_P=0.
\]
\end{proof}

\subsection{Defect characteristic classes}

\begin{definition}
The cohomology class
\[
[\omega_P]
\in
H^{2k}_{\mathrm{CE}}(\mathfrak g)
\]
is called the \emph{characteristic class associated with the equivariance
defect of \(D\)}.
\end{definition}

\begin{remark}
This class measures the obstruction to infinitesimal equivariance of the
operator \(D\).
\end{remark}

\subsection{Naturality}

Let
\[
f:X\to B^{\mathrm{sph}}(G)
\]
be a classifying map, and let
\[
P_f=f^*E^{\mathrm{sph}}(G)
\]
be the associated principal bundle.

\begin{proposition}
The defect construction pulls back naturally to
\[
P_f.
\]
\end{proposition}

\begin{proof}
The pulled-back bundle inherits the pulled-back Clifford module and the
pulled-back operator
\[
D_f.
\]
The infinitesimal action of \(\mathfrak g\) on
\[
P_f
\]
induces the corresponding defect
\[
\mathcal D_{D_f}.
\]
All constructions defining
\[
\mathcal F_D
\]
and
\[
\omega_P
\]
commute with pullback.
\end{proof}

\begin{proposition}
The defect characteristic classes are functorial under pullback of principal
bundles.
\end{proposition}

\subsection{Homotopy invariance}

\begin{proposition}
The defect characteristic classes depend only on the homotopy class of the
classifying map.
\end{proposition}

\begin{proof}
Characteristic classes obtained through Chern--Weil-type constructions are
natural under pullback. Since homotopic maps induce isomorphic pullback
principal bundles, the corresponding defect classes coincide.
\end{proof}

\subsection{Vanishing criterion}

\begin{proposition}
If
\[
\mathcal D_D=0,
\]
then all defect characteristic classes vanish.
\end{proposition}

\begin{proof}
If
\[
\mathcal D_D=0,
\]
then
\[
\mathcal F_D=0.
\]
Therefore every invariant polynomial evaluated on
\[
\mathcal F_D
\]
vanishes identically, and so do the associated cohomology classes.
\end{proof}

\begin{remark}
The converse is generally false: the characteristic classes detect only the
cohomological part of the defect.
\end{remark}

\subsection{Interpretation}

\begin{remark}
The equivariance defect behaves as a secondary curvature measuring the failure
of a geometric operator to descend from
\[
E^{\mathrm{sph}}(G)
\]
to
\[
B^{\mathrm{sph}}(G).
\]
The associated characteristic classes are therefore secondary obstruction
classes attached not to a connection, but to the failure of equivariance of an
operator.
\end{remark}

\begin{remark}
The construction is formally analogous to Chern--Weil theory, but it naturally
belongs to Lie algebra cohomology rather than ordinary de Rham cohomology on
the quotient space.
\end{remark}
\section{Defect classes, Schwinger cocycles and lifting gerbes}
\label{sec:defect-schwinger-gerbe}

We now relate equivariance defects to central extensions and lifting gerbes.
The construction requires additional analytic hypotheses. In particular, one
must work in a restricted Hilbert setting where the Schwinger cocycle is
well-defined.

\subsection{Restricted Hilbert setting}

Let
\[
H=H_+\oplus H_-
\]
be a polarized Hilbert space, and let
\[
\mathrm{GL}_{\mathrm{res}}(H)
\]
be the restricted general linear group, i.e. the group of invertible operators
whose off-diagonal blocks are Hilbert--Schmidt with respect to the
polarization \cite{PressleySegal1986}.

Its Lie algebra is
\[
\mathfrak{gl}_{\mathrm{res}}(H)
=
\left\{
A\in\mathcal B(H)
\;\middle|\;
A_{+-},A_{-+}\text{ are Hilbert--Schmidt}
\right\}.
\]

Let
\[
\epsilon
=
\begin{pmatrix}
1 & 0\\
0 & -1
\end{pmatrix}
\]
be the grading operator associated with the polarization.

\subsection{The Schwinger cocycle}

\begin{definition}
The Schwinger cocycle on
\[
\mathfrak{gl}_{\mathrm{res}}(H)
\]
is the bilinear map
\[
c_{\mathrm{Sch}}:
\mathfrak{gl}_{\mathrm{res}}(H)\times
\mathfrak{gl}_{\mathrm{res}}(H)
\longrightarrow
\mathbb C
\]
defined by
\[
c_{\mathrm{Sch}}(A,B)
=
\operatorname{Tr}
\left(
A_{-+}B_{+-}
-
B_{-+}A_{+-}
\right).
\]
Equivalently,
\[
c_{\mathrm{Sch}}(A,B)
=
\frac14\operatorname{Tr}
\left(
\epsilon[\epsilon,A][\epsilon,B]
\right),
\]
whenever the latter expression is interpreted in the standard restricted
sense.
\end{definition}

\begin{proposition}
The map
\[
c_{\mathrm{Sch}}
\]
is a Lie algebra \(2\)-cocycle on
\[
\mathfrak{gl}_{\mathrm{res}}(H).
\]
\end{proposition}

\begin{proof}
The off-diagonal blocks of \(A\) and \(B\) are Hilbert--Schmidt, hence the
products
\[
A_{-+}B_{+-},
\qquad
B_{-+}A_{+-}
\]
are trace-class. Thus \(c_{\mathrm{Sch}}(A,B)\) is well-defined.

Skew-symmetry is immediate from the formula.

The cocycle identity follows from the standard restricted-trace computation:
for
\[
A,B,C\in\mathfrak{gl}_{\mathrm{res}}(H),
\]
one checks that
\[
c_{\mathrm{Sch}}([A,B],C)
+
c_{\mathrm{Sch}}([B,C],A)
+
c_{\mathrm{Sch}}([C,A],B)
=
0.
\]
This is the usual Schwinger cocycle on the restricted Lie algebra; see
\cite{PressleySegal1986}.
\end{proof}

\subsection{Pullback by a restricted infinitesimal action}

Assume now that the \(G\)-action on the relevant Hilbert module is represented
infinitesimally by a Lie algebra morphism
\[
\rho_*:\mathfrak g\longrightarrow \mathfrak{gl}_{\mathrm{res}}(H).
\]

\begin{definition}
The Schwinger cocycle associated with the restricted action is
\[
c_\rho(\xi,\eta)
:=
c_{\mathrm{Sch}}(\rho_*(\xi),\rho_*(\eta)).
\]
\end{definition}

\begin{proposition}
The bilinear map
\[
c_\rho:\mathfrak g\times\mathfrak g\to\mathbb C
\]
is a Lie algebra \(2\)-cocycle.
\end{proposition}

\begin{proof}
Since \(\rho_*\) is a Lie algebra morphism,
\[
\rho_*([\xi,\eta])
=
[\rho_*(\xi),\rho_*(\eta)].
\]
The cocycle identity for \(c_\rho\) is therefore the pullback of the cocycle
identity for \(c_{\mathrm{Sch}}\).
\end{proof}

\subsection{Defect-induced Schwinger cocycle}

We now incorporate the equivariance defect.

Let
\[
D
\]
be a Dirac-type operator, and suppose that its infinitesimal equivariance
defect is defined by
\[
\mathcal D_D(\xi)
=
[\mathcal L_{\xi^\#},D].
\]

\textbf{Assumption.}
We assume that the defect operators define a map
\[
\delta_D:\mathfrak g\longrightarrow \mathfrak{gl}_{\mathrm{res}}(H),
\qquad
\xi\mapsto \mathcal D_D(\xi),
\]
and that \(\delta_D\) is a Lie algebra morphism, or more generally that its
failure to be a morphism is controlled by trace-class terms for which the
Schwinger cocycle remains defined.


\begin{definition}
Under this assumption, the \emph{defect Schwinger cocycle} is
\[
c_D(\xi,\eta)
=
c_{\mathrm{Sch}}(\delta_D(\xi),\delta_D(\eta)).
\]
Explicitly,
\[
c_D(\xi,\eta)
=
\operatorname{Tr}
\left(
\delta_D(\xi)_{-+}\delta_D(\eta)_{+-}
-
\delta_D(\eta)_{-+}\delta_D(\xi)_{+-}
\right).
\]
\end{definition}

\begin{proposition}
If \(\delta_D\) is a Lie algebra morphism into
\(\mathfrak{gl}_{\mathrm{res}}(H)\), then
\[
c_D
\]
is a Lie algebra \(2\)-cocycle on \(\mathfrak g\).
\end{proposition}

\begin{proof}
This is again the pullback of the Schwinger cocycle:
\[
c_D=\delta_D^*c_{\mathrm{Sch}}.
\]
Since \(c_{\mathrm{Sch}}\) is closed in Lie algebra cohomology and
\(\delta_D\) is a Lie algebra morphism, \(c_D\) is closed.
\end{proof}

\begin{remark}
The assumption that \(\delta_D\) be a Lie algebra morphism is restrictive. If
it is not satisfied, then \(c_D\) should be regarded first as a bilinear
Schwinger-type functional. Its failure to be closed is measured by the
Chevalley--Eilenberg differential
\[
d_{\mathrm{CE}}c_D.
\]
Thus the cocycle property is an additional analytic and algebraic condition,
not an automatic consequence of the definition of the defect.
\end{remark}

\subsection{Central extension}

Assume that
\[
c_D
\]
is a Lie algebra \(2\)-cocycle with values in \(i\mathbb R\) or \(\mathbb R\),
depending on the chosen convention.

\begin{proposition}
The cocycle \(c_D\) defines a central extension of Lie algebras
\[
0
\longrightarrow
\mathbb R
\longrightarrow
\widehat{\mathfrak g}_D
\longrightarrow
\mathfrak g
\longrightarrow
0.
\]
\end{proposition}

\begin{proof}
Define
\[
\widehat{\mathfrak g}_D
=
\mathfrak g\oplus\mathbb R
\]
with bracket
\[
[(\xi,a),(\eta,b)]
=
([\xi,\eta],c_D(\xi,\eta)).
\]
The Jacobi identity for this bracket is equivalent to the cocycle identity for
\(c_D\). The second factor is central, so the sequence is a central extension.
\end{proof}

\subsection{Integration to a group extension}

\textbf{assumption}
We assume that the Lie algebra extension defined by \(c_D\) integrates to a
central extension of diffeological Lie groups
\[
1
\longrightarrow
U(1)
\longrightarrow
\widehat G_D
\longrightarrow
G
\longrightarrow
1.
\]


\begin{remark}
This integration is not automatic. It requires an integrality condition on the
periods of the Lie algebra cocycle and suitable regularity assumptions on
\(G\).
\end{remark}

\subsection{The lifting gerbe}

Let
\[
E^{\mathrm{sph}}(G)\to B^{\mathrm{sph}}(G)
\]
be the universal principal \(G\)-bundle.

\begin{definition}
The \emph{Schwinger lifting gerbe} associated with \(D\) is the lifting gerbe
for the central extension
\[
1
\to
U(1)
\to
\widehat G_D
\to
G
\to
1
\]
applied to the universal principal \(G\)-bundle
\[
E^{\mathrm{sph}}(G)\to B^{\mathrm{sph}}(G).
\]
\end{definition}

\begin{proposition}
The Schwinger lifting gerbe defines a class
\[
\mathrm{DD}(D)
\in
H^3(B^{\mathrm{sph}}(G);\mathbb Z).
\]
\end{proposition}

\begin{proof}
A central extension of \(G\) by \(U(1)\) defines a lifting problem for any
principal \(G\)-bundle. The obstruction to lifting a principal \(G\)-bundle to
a principal \(\widehat G_D\)-bundle is represented by a \(U(1)\)-banded gerbe.
Its Dixmier--Douady class lies in degree-three integral cohomology. Applying
this construction to the universal spherical bundle gives the stated class.
\end{proof}

\subsection{Pullback by classifying maps}

Let
\[
f:X\to B^{\mathrm{sph}}(G)
\]
be a classifying map, and let
\[
P=f^*E^{\mathrm{sph}}(G)
\]
be the associated principal \(G\)-bundle.

\begin{theorem}
The pullback class
\[
f^*\mathrm{DD}(D)
\in
H^3(X;\mathbb Z)
\]
is the Dixmier--Douady class of the lifting gerbe of \(P\) associated with the
central extension
\[
1\to U(1)\to \widehat G_D\to G\to 1.
\]
\end{theorem}

\begin{proof}
The lifting gerbe construction is functorial under pullback of principal
bundles. Since
\[
P=f^*E^{\mathrm{sph}}(G),
\]
the lifting gerbe of \(P\) is the pullback of the universal lifting gerbe on
\(B^{\mathrm{sph}}(G)\). Therefore its Dixmier--Douady class is
\[
f^*\mathrm{DD}(D).
\]
\end{proof}

\subsection{Interpretation}

\begin{remark}
The class
\[
\mathrm{DD}(D)
\]
should be interpreted as the higher obstruction associated with the
Schwinger-type defect of \(D\). It measures the obstruction to lifting the
\(G\)-symmetry to the centrally extended symmetry group
\[
\widehat G_D.
\]
\end{remark}

\begin{remark}
In field-theoretic terminology, this is an anomaly class: the Schwinger term
is the infinitesimal cocycle, while the Dixmier--Douady class is its global
gerbe-theoretic counterpart.
\end{remark}

\subsection{Summary}

Under the restricted Hilbert and integrability assumptions above, the
equivariance defect of a Dirac-type operator gives rise to:
\[
\text{defect}
\quad\Rightarrow\quad
\text{Schwinger cocycle}
\quad\Rightarrow\quad
\text{central extension}
\quad\Rightarrow\quad
\text{lifting gerbe}
\quad\Rightarrow\quad
\text{Dixmier--Douady class}.
\]
\section{An explicit example: loop groups and the Schwinger cocycle}
\label{sec:loop-example}

We now describe an explicit model in which the defect class can be computed.
This example is classical and provides the prototype for the construction above.

\subsection{The polarized Hilbert space}

Let
\[
H=L^2(S^1,\mathbb C^N)
\]
and consider the Fourier polarization
\[
H=H_+\oplus H_-,
\]
where
\[
H_+
=
\overline{\operatorname{span}}\{z^n e_a \mid n\geq0,\ 1\leq a\leq N\},
\]
and
\[
H_-
=
\overline{\operatorname{span}}\{z^n e_a \mid n<0,\ 1\leq a\leq N\}.
\]

Let
\[
D=-i\frac{d}{d\theta}
\]
be the standard Dirac operator on the circle. Its positive and negative spectral subspaces are precisely \(H_+\) and \(H_-\), up to the convention for the zero mode.

\subsection{The loop algebra action}

Let \(K\) be a compact Lie group with Lie algebra \(\mathfrak k\), and let
\[
L\mathfrak k=C^\infty(S^1,\mathfrak k)
\]
be its loop algebra.

Each \(X\in L\mathfrak k\) acts on \(H\) by pointwise multiplication:
\[
(M_X\psi)(\theta)=X(\theta)\psi(\theta).
\]

With respect to the polarization \(H=H_+\oplus H_-\), the operator \(M_X\) has off-diagonal blocks of Hilbert--Schmidt type. Hence
\[
M_X\in \mathfrak{gl}_{\mathrm{res}}(H).
\]

\subsection{The Schwinger cocycle}

\begin{definition}
The Schwinger cocycle on \(L\mathfrak k\) is defined by
\[
c(X,Y)
=
\operatorname{Tr}\bigl(
(M_X)_{-+}(M_Y)_{+-}
-
(M_Y)_{-+}(M_X)_{+-}
\bigr),
\]
where the block decomposition is taken with respect to
\[
H=H_+\oplus H_-.
\]
\end{definition}

\begin{theorem}
The cocycle \(c\) is given explicitly by
\[
c(X,Y)
=
\frac{1}{2\pi i}
\int_{S^1}
\operatorname{tr}\bigl(X\,dY\bigr).
\]
Equivalently, if \(z=e^{i\theta}\), then
\[
c(X,Y)
=
\frac{1}{2\pi i}
\int_{|z|=1}
\operatorname{tr}\left(X(z)\frac{dY}{dz}(z)\right)\,dz.
\]
\end{theorem}

\begin{proof}
Write the Fourier expansions
\[
X(z)=\sum_{m\in\mathbb Z}X_m z^m,
\qquad
Y(z)=\sum_{n\in\mathbb Z}Y_n z^n.
\]

Multiplication by \(X\) sends the Fourier mode \(z^r\) to
\[
\sum_m X_m z^{m+r}.
\]

The block \((M_X)_{-+}\) consists of terms sending \(H_+\) to \(H_-\), hence those for which
\[
r\geq0,\qquad m+r<0.
\]
Similarly, \((M_Y)_{+-}\) consists of terms sending \(H_-\) to \(H_+\).

A direct trace computation gives
\[
\operatorname{Tr}\bigl((M_X)_{-+}(M_Y)_{+-}\bigr)
-
\operatorname{Tr}\bigl((M_Y)_{-+}(M_X)_{+-}\bigr)
=
\sum_{m\in\mathbb Z} m\, \operatorname{tr}(X_{-m}Y_m).
\]

On the other hand,
\[
dY
=
\sum_m mY_m z^{m-1}dz.
\]
Thus
\[
\operatorname{tr}(X\,dY)
=
\sum_{m,n}
n\,\operatorname{tr}(X_mY_n)z^{m+n-1}dz.
\]
The residue picks out the terms \(m+n=0\), so
\[
\frac{1}{2\pi i}
\int_{|z|=1}
\operatorname{tr}(X\,dY)
=
\sum_n n\,\operatorname{tr}(X_{-n}Y_n).
\]
This is exactly the Schwinger cocycle.
\end{proof}

\subsection{Central extension}

\begin{proposition}
The cocycle \(c\) defines a central extension
\[
0\longrightarrow \mathbb C
\longrightarrow
\widehat{L\mathfrak k}
\longrightarrow
L\mathfrak k
\longrightarrow 0,
\]
with bracket
\[
[(X,a),(Y,b)]
=
([X,Y],c(X,Y)).
\]
\end{proposition}

\begin{proof}
The bilinear form \(c\) is skew-symmetric and satisfies the Lie algebra cocycle identity. Hence it defines a central extension of \(L\mathfrak k\).
\end{proof}

\subsection{Relation with the defect of the Dirac operator}

Let \(D=-i\frac{d}{d\theta}\). For \(X\in L\mathfrak k\), compute
\[
[D,M_X]
=
-iM_{X'},
\]
where
\[
X'=\frac{dX}{d\theta}.
\]

Thus the failure of multiplication operators to commute with the Dirac operator is measured by differentiation of the loop.

\begin{proposition}
The Schwinger cocycle can be expressed in terms of the Dirac defect by
\[
c(X,Y)
=
\operatorname{Tr}_{\mathrm{res}}
\bigl(
M_X[D,M_Y]
\bigr).
\]
\end{proposition}

\begin{proof}
Since
\[
[D,M_Y]=-iM_{Y'},
\]
we obtain
\[
\operatorname{Tr}_{\mathrm{res}}
\bigl(M_X[D,M_Y]\bigr)
=
-i\,\operatorname{Tr}_{\mathrm{res}}(M_XM_{Y'}).
\]
The restricted trace extracts the residue contribution of the commutator with the polarization, and this is precisely
\[
\frac{1}{2\pi i}
\int_{S^1}\operatorname{tr}(X\,dY).
\]
Hence the two expressions agree up to the conventional normalization.
\end{proof}

\begin{remark}
This formula shows explicitly that the Schwinger term is the trace-level shadow of the failure of the loop algebra action to commute with the Dirac operator.
\end{remark}

\subsection{Gerbe class}

The central extension of \(L\mathfrak k\) integrates, under the usual integrality assumptions, to a central extension of the loop group
\[
1\longrightarrow U(1)
\longrightarrow
\widehat{LK}
\longrightarrow
LK
\longrightarrow
1.
\]

This central extension defines a lifting gerbe over the classifying space
\[
B LK.
\]

\begin{definition}
The Dixmier--Douady class associated with the Schwinger cocycle is
\[
\operatorname{DD}_{\mathrm{Sch}}
\in
H^3(BLK,\mathbb Z).
\]
\end{definition}

\begin{proposition}
For a principal \(LK\)-bundle
\[
P\to X
\]
classified by
\[
f:X\to BLK,
\]
the obstruction to lifting \(P\) to a principal \(\widehat{LK}\)-bundle is
\[
f^*\operatorname{DD}_{\mathrm{Sch}}
\in
H^3(X,\mathbb Z).
\]
\end{proposition}

\begin{proof}
The central extension
\[
1\to U(1)\to \widehat{LK}\to LK\to1
\]
defines a lifting gerbe over \(BLK\). Pulling this gerbe back along the classifying map \(f\) gives the lifting gerbe of \(P\). Its Dixmier--Douady class is therefore
\[
f^*\operatorname{DD}_{\mathrm{Sch}}.
\]
The vanishing of this class is equivalent to the existence of a lift of \(P\) to a \(\widehat{LK}\)-bundle.
\end{proof}

\subsection{Interpretation for spherical Milnor spaces}

In the setting of spherical Milnor classifying spaces, one may take
\[
G=LK.
\]
Then the universal spherical bundle
\[
E^{\mathrm{sph}}(LK)\to B^{\mathrm{sph}}(LK)
\]
carries the same Schwinger defect mechanism.

The Dirac operator on the circle provides the polarization, the loop algebra acts by restricted operators, and the defect of equivariance produces the cocycle
\[
c(X,Y)
=
\frac{1}{2\pi i}
\int_{S^1}
\operatorname{tr}(X\,dY).
\]

Thus the universal defect class on
\[
B^{\mathrm{sph}}(LK)
\]
pulls back to the usual Schwinger--Dixmier--Douady class on any space classified by a map
\[
X\to B^{\mathrm{sph}}(LK).
\]

\begin{remark}
This example shows that the defect characteristic class introduced above is not merely formal: in the classical loop group case, it reproduces the standard Schwinger cocycle and the associated Dixmier--Douady obstruction.
\end{remark}

\section{Current groups, spin manifolds and non-spin twists}
\label{sec:current-groups-spin}

We now extend the loop group example to current groups over compact manifolds. This provides a higher-dimensional source of Schwinger-type cocycles and defect gerbes.

\subsection{Current groups over a spin manifold}

Let \(M\) be a compact oriented Riemannian spin manifold of dimension \(m\), and let
\[
K
\]
be a compact Lie group with Lie algebra \(\mathfrak k\).

The current group is
\[
\mathcal G_M:=C^\infty(M,K),
\]
and its Lie algebra is
\[
\mathfrak g_M:=C^\infty(M,\mathfrak k).
\]

Let
\[
S\to M
\]
be the spinor bundle, and let
\[
E:=M\times V
\]
be the trivial vector bundle associated with a finite-dimensional unitary representation
\[
K\to U(V).
\]

The Hilbert space is
\[
\mathcal H=L^2(M,S\otimes V).
\]

The Dirac operator is
\[
D_M:\Gamma(S\otimes V)\longrightarrow \Gamma(S\otimes V).
\]

The spectral decomposition of \(D_M\) gives a polarization
\[
\mathcal H=\mathcal H_+\oplus \mathcal H_-,
\]
where \(\mathcal H_+\) and \(\mathcal H_-\) are the positive and negative spectral subspaces of \(D_M\), with a finite-dimensional convention for the kernel.

\subsection{Action of the current algebra}

An element
\[
X\in C^\infty(M,\mathfrak k)
\]
acts on \(\mathcal H\) by pointwise multiplication:
\[
(M_X\psi)(p)=X(p)\psi(p).
\]

The commutator with the Dirac operator is
\[
[D_M,M_X]=c(dX),
\]
where \(c\) denotes Clifford multiplication.

\begin{proposition}
The current algebra acts by restricted operators with respect to the spectral polarization of \(D_M\), under the usual Schatten summability condition determined by \(\dim M\).
\end{proposition}

\begin{proof}
The operator \(M_X\) is a zeroth-order pseudodifferential operator. Its commutator with the sign operator
\[
F:=D_M|D_M|^{-1}
\]
is of order \(-1\). On an \(m\)-dimensional compact manifold, a pseudodifferential operator of order \(-1\) belongs to Schatten classes \(L^p\) for \(p>m\). Thus the off-diagonal components of \(M_X\) with respect to the polarization are compact and belong to the appropriate restricted ideal. In dimension \(1\), this is Hilbert--Schmidt; in higher dimensions, one must use the \(p\)-restricted Grassmannian.
\end{proof}

\begin{remark}
Thus the correct higher-dimensional analogue of the loop group construction is not always the Hilbert--Schmidt restricted Grassmannian, but rather a Schatten-restricted Grassmannian adapted to the dimension of \(M\).
\end{remark}

\subsection{Schwinger cocycle in the spin case}

The defect of equivariance of the Dirac operator is
\[
\delta_D(X):=[D_M,M_X]=c(dX).
\]

The corresponding Schwinger-type cocycle is obtained by regularized trace expressions involving the polarization \(F\). Formally, one writes
\[
c_D(X,Y)
=
\operatorname{Tr}_{\mathrm{reg}}
\left(
M_X[F,M_Y]
\right).
\]

In dimension \(1\), this reduces to
\[
c_D(X,Y)
=
\frac{1}{2\pi i}\int_{S^1}\operatorname{tr}(X\,dY).
\]

In higher odd dimensions, the local index formula gives higher Schwinger cocycles of the form
\[
c_D(X_0,\ldots,X_m)
=
C_m\int_M
\operatorname{tr}
\left(
X_0\,dX_1\cdots dX_m
\right),
\]
up to normalization constants and characteristic factors determined by the Dirac operator.

\begin{proposition}
The higher Schwinger cocycle associated with \(D_M\) defines a Lie algebra cohomology class of the current algebra \(C^\infty(M,\mathfrak k)\).
\end{proposition}

\begin{proof}
The expression is obtained from the transgression of the Chern character of the Dirac spectral triple. The cocycle identity follows from the closedness of the corresponding local index density and the cyclicity of the regularized trace. Equivalently, it is the current algebra cocycle determined by the descent equations for the Chern character.
\end{proof}

\subsection{Central extensions and higher gerbes}

For \(M=S^1\), the Schwinger cocycle integrates to a central extension
\[
1\to U(1)\to \widehat{LK}\to LK\to 1.
\]

For higher-dimensional \(M\), the corresponding object is no longer, in general, an ordinary central extension by \(U(1)\). Instead, one obtains higher extensions or gerbes over the current group.

\begin{definition}
The defect gerbe of the current group \(C^\infty(M,K)\) is the gerbe whose Dixmier--Douady-type class is represented by the higher Schwinger cocycle associated with \(D_M\).
\end{definition}

\begin{remark}
In dimension \(1\), this reduces to the usual lifting gerbe associated with the central extension of the loop group. In higher dimensions, it is naturally a higher gerbe or a higher group extension.
\end{remark}

\subsection{Pullback to spherical Milnor classifying spaces}

Take
\[
G=C^\infty(M,K).
\]

The spherical Milnor classifying space
\[
B^{\mathrm{sph}}(G)
\]
carries the universal principal current group bundle.

The Dirac operator \(D_M\) produces a universal defect class
\[
\beta_D^{\mathrm{univ}}
\]
on
\[
B^{\mathrm{sph}}(C^\infty(M,K)).
\]

For any classifying map
\[
f:X\to B^{\mathrm{sph}}(C^\infty(M,K)),
\]
the pullback
\[
f^*\beta_D^{\mathrm{univ}}
\]
is the defect class of the associated current group bundle over \(X\).

\begin{theorem}
In the spin case, the defect class associated with the current group
\[
C^\infty(M,K)
\]
is determined by the Dirac operator on \(M\) and is functorial under pullback by classifying maps.
\end{theorem}

\begin{proof}
The Dirac operator defines the polarization and hence the restricted current algebra representation. The Schwinger cocycle is constructed functorially from the commutators
\[
[D_M,M_X].
\]
Therefore the resulting gerbe or higher gerbe class is natural with respect to pullback of principal current group bundles. Since classifying maps pull back the universal bundle, the associated defect class is the pullback of the universal one.
\end{proof}

\subsection{The non-spin case}

We now consider the case where \(M\) is not spin.

If \(M\) is not spin, there is no globally defined spinor bundle \(S\), and therefore no ordinary Dirac operator on spinors. Consequently, the construction above cannot be applied directly.

There are three natural replacements.

\subsubsection{The \(\mathrm{Spin}^c\) case}

Assume that \(M\) admits a \(\mathrm{Spin}^c\)-structure. Then there exists a complex spinor bundle
\[
S^c\to M
\]
and a \(\mathrm{Spin}^c\)-Dirac operator
\[
D_M^c:\Gamma(S^c\otimes V)\to \Gamma(S^c\otimes V).
\]

The current algebra acts again by multiplication, and the defect is
\[
\delta_{D^c}(X)=[D_M^c,M_X]=c(dX).
\]

\begin{proposition}
For a \(\mathrm{Spin}^c\)-manifold \(M\), the current group defect construction goes through with \(D_M\) replaced by \(D_M^c\).
\end{proposition}

\begin{proof}
The proof is identical to the spin case. The \(\mathrm{Spin}^c\)-Dirac operator is a first-order elliptic operator with the same principal symbol given by Clifford multiplication. Hence the commutator with multiplication by a current is again
\[
[D_M^c,M_X]=c(dX).
\]
The spectral polarization and the corresponding Schwinger cocycles are therefore defined in the same way.
\end{proof}

\begin{remark}
The resulting defect class now depends on the chosen \(\mathrm{Spin}^c\)-structure and on the determinant line bundle of that structure.
\end{remark}

\subsubsection{The twisted spin case}

If \(M\) is not spin but has a spin lifting gerbe, one may replace the spinor bundle by a twisted spinor module.

The obstruction to a spin structure is
\[
w_2(M)\in H^2(M;\mathbb Z_2).
\]

Its integral lift, when present through the Bockstein homomorphism, gives the obstruction
\[
W_3(M)\in H^3(M;\mathbb Z).
\]

\begin{definition}
A twisted spinor bundle is a module over the Clifford algebra bundle defined locally on spin charts, with transition functions twisted by the spin lifting gerbe.
\end{definition}

In this case, the Dirac operator is not an ordinary global operator on a vector bundle, but a twisted Dirac operator acting on sections of a gerbe module.

\begin{proposition}
For a non-spin manifold admitting a twisted spinor module, the defect construction defines a gerbe-twisted Schwinger class.
\end{proposition}

\begin{proof}
Locally, on spin charts, the construction is identical to the spin case. On overlaps, local spinor bundles differ by the transition data of the spin lifting gerbe. The local Schwinger cocycles therefore glue not as ordinary differential forms, but as cocycles with coefficients twisted by the gerbe. The resulting global object is a gerbe-twisted defect class.
\end{proof}

\subsubsection{The de Rham-type replacement}

If no spin or \(\mathrm{Spin}^c\)-structure is chosen, one can still use the de Rham operator
\[
D_{\mathrm{dR}}=d+d^*
\]
acting on differential forms.

The current algebra acts by multiplication through a representation of \(K\) on an auxiliary vector bundle. Then
\[
[D_{\mathrm{dR}},M_X]
\]
is again a first-order defect controlled by \(dX\).

\begin{proposition}
The de Rham operator produces a canonical defect construction on any compact oriented Riemannian manifold.
\end{proposition}

\begin{proof}
The de Rham operator exists without a spin structure. It is elliptic and first order. Multiplication by a smooth current has commutator
\[
[D_{\mathrm{dR}},M_X]
\]
given by exterior multiplication and contraction with \(dX\), depending on the representation. Hence it produces a well-defined operator defect and an associated regularized trace cocycle, whenever the corresponding Schatten conditions are satisfied.
\end{proof}

\begin{remark}
This replacement is less spin-geometric, but it is canonical on any oriented Riemannian manifold.
\end{remark}

\subsection{Comparison of the three cases}

The situation can be summarized as follows.

\[
\begin{array}{c|c|c|c}
\text{Structure on }M & \text{Operator} & \text{Defect class} & \text{Nature} \\
\hline
\text{spin} & D_M & \beta_D & \text{ordinary Schwinger/gerbe class} \\
\mathrm{Spin}^c & D_M^c & \beta_{D^c} & \text{depends on determinant line} \\
\text{twisted spin} & D_{\mathrm{tw}} & \beta_{\mathrm{tw}} & \text{gerbe-twisted class} \\
\text{oriented Riemannian} & d+d^* & \beta_{\mathrm{dR}} & \text{de Rham defect class}
\end{array}
\]

\begin{remark}
Thus the spin case gives the cleanest analogue of the loop group Schwinger cocycle. The non-spin case does not destroy the construction, but shifts it from ordinary spinorial geometry to either \(\mathrm{Spin}^c\)-geometry, gerbe-twisted spin geometry, or de Rham-type geometry.
\end{remark}

\subsection{Conclusion}

Current groups over spin manifolds provide the natural higher-dimensional generalization of the loop group example. The defect of equivariance of the Dirac operator produces Schwinger-type cocycles and higher gerbe classes.

For non-spin manifolds, the construction persists only after replacing the spin Dirac operator by an appropriate substitute: a \(\mathrm{Spin}^c\)-Dirac operator, a gerbe-twisted Dirac operator, or the de Rham operator.

\vskip 12pt

\paragraph{Conflict of interest statement:} The author declares no conflict of interest.

\vskip 12pt

\paragraph{\bf Acknowledgements:} J.-P.M  thanks the France 2030 framework programme Centre Henri Lebesgue ANR-11-LABX-0020-01 
for creating an attractive mathematical environment.

\vskip 12pt

\paragraph{\bf Declaration of generative AI and AI-assisted technologies in the writing process}

During the preparation of this work the author used ChatGPT and Mistral AI in order to smoothen the expression in English. After using this tool/service, the author reviewed and edited the content as needed and takes full responsibility for the content of the publication.

\begin{thebibliography}{99}

\bibitem{AtiyahSegal2004}
M.~F.~Atiyah and G.~Segal,
\newblock Twisted K-theory,
\newblock {\em Ukrainian Math. Bull.} 1 (2004), no.~3, 291--334.




\bibitem{BaezLauda2004}
J.~C.~Baez and A.~D.~Lauda,
\newblock Higher-dimensional algebra V: 2-groups,
\newblock {\em Theory Appl. Categ.} 12 (2004), 423--491.

\bibitem{Breen1994}
L.~Breen,
\newblock {\em On the Classification of 2-Gerbes and 2-Stacks},
\newblock Ast\'erisque, vol.~225,
\newblock Soci\'et\'e Math\'ematique de France, Paris, 1994.

\bibitem{Breen2008}
L.~Breen,
\newblock Differential geometry of gerbes and differential forms,
\newblock arXiv:0802.1833.

\bibitem{BreenMessing2005}
L.~Breen and W.~Messing,
\newblock Differential geometry of gerbes,
\newblock {\em Adv. Math.} 198 (2005), no.~2, 732--846.

\bibitem{BrownHiggins1981}
R.~Brown and P.~J.~Higgins,
\newblock The algebra of cubes,
\newblock {\em J. Pure Appl. Algebra} 21 (1981), no.~3, 233--260.

\bibitem{BrownHigginsSivera2011}
R.~Brown, P.~J.~Higgins, and R.~Sivera,
\newblock {\em Nonabelian Algebraic Topology},
\newblock EMS Tracts in Mathematics, vol.~15,
\newblock European Mathematical Society, Z\"urich, 2011.

\bibitem{CareyMickelssonMurray1997}
A.~L.~Carey, J.~Mickelsson, and M.~K.~Murray,
\newblock Index theory, gerbes, and Hamiltonian quantization,
\newblock {\em Comm. Math. Phys.} 183 (1997), no.~3, 707--722.

\bibitem{ChristensenSinnamonWu2014}
J.~D.~Christensen, G.~Sinnamon, and E.~Wu,
\newblock The $D$-topology for diffeological spaces,
\newblock {\em Pacific J. Math.} 272 (2014), no.~1, 87--110.

\bibitem{ChristensenWu2014Homotopy}
J.~D.~Christensen and E.~Wu,
\newblock The homotopy theory of diffeological spaces,
\newblock {\em New York J. Math.} 20 (2014), 1269--1303.

\bibitem{ChristensenWu2016Tangent}
J.~D.~Christensen and E.~Wu,
\newblock Tangent spaces and tangent bundles for diffeological spaces,
\newblock {\em Cah. Topol. G\'eom. Diff\'er. Cat\'eg.} 57 (2016), no.~1, 3--50.

\bibitem{ChristensenWu2017Filtered}
J.~D.~Christensen and E.~Wu,
\newblock Tangent spaces of bundles and of filtered diffeological spaces,
\newblock {\em Proc. Amer. Math. Soc.} 145 (2017), no.~5, 2255--2270.

\bibitem{ChristensenWu2019Vector}
J.~D.~Christensen and E.~Wu,
\newblock Diffeological vector spaces,
\newblock {\em Pacific J. Math.} 303 (2019), no.~1, 73--92.

\bibitem{ChristensenWu2021Classifying}
J.~D.~Christensen and E.~Wu,
\newblock Smooth classifying spaces,
\newblock {\em Israel J. Math.} 241 (2021), no.~2, 911--954.

\bibitem{ChristensenWu2022Exterior}
J.~D.~Christensen and E.~Wu,
\newblock Exterior bundles in diffeology,
\newblock {\em Israel J. Math.} 251 (2022), 537--581.



\bibitem{Dold1963}
A.~Dold,
\newblock Partitions of unity in the theory of fibrations,
\newblock {\em Ann. of Math.} (2) 78 (1963), 223--255.

\bibitem{DonovanKaroubi1970}
P.~Donovan and M.~Karoubi,
\newblock Graded Brauer groups and K-theory with local coefficients,
\newblock {\em Inst. Hautes \'Etudes Sci. Publ. Math.} 38 (1970), 5--25.

\bibitem{Dupont1978}
J.~L.~Dupont,
\newblock {\em Curvature and Characteristic Classes},
\newblock Lecture Notes in Mathematics, vol.~640,
\newblock Springer, Berlin--New York, 1978.

\bibitem{FiorenzaSchreiberStasheff2007}
D.~Fiorenza, U.~Schreiber, and J.~Stasheff,
\newblock \v{C}ech cocycles for differential characteristic classes:
\newblock an $\infty$-Lie theoretic construction,
\newblock {\em Adv. Theor. Math. Phys.} 16 (2012), no.~1, 149--250.

\bibitem{Freed1995}
D.~S.~Freed,
\newblock Classical Chern--Simons theory. Part 1,
\newblock {\em Adv. Math.} 113 (1995), no.~2, 237--303.

\bibitem{Giraud1971}
J.~Giraud,
\newblock {\em Cohomologie non ab\'elienne},
\newblock Grundlehren der mathematischen Wissenschaften, vol.~179,
\newblock Springer-Verlag, Berlin--New York, 1971.

	\bibitem{GMW2023} Goldammer, N.; Magnot, J-P.; Welker, K.; On diffeologies from infinite dimensional geometry to PDE constrained optimization \emph{Contemporary Mathematics} {\bf 794} (2024), 1--48.

\bibitem{Hector1995}
G.~Hector,
\newblock G\'eom\'etrie et topologie des espaces diff\'eologiques,
\newblock in {\em Analysis and Geometry in Foliated Manifolds},
\newblock World Scientific, Singapore, 1995, pp.~55--80.

\bibitem{IglesiasLaffineurZemmour2020}
P.~Iglesias-Zemmour and F.~Laffineur,
\newblock Noncommutative geometry and diffeology: the irrational torus,
\newblock {\em J. Noncommut. Geom.} 14 (2020), no.~4, 1249--1282.

\bibitem{IglesiasZemmour2013}
P.~Iglesias-Zemmour,
\newblock {\em Diffeology},
\newblock Mathematical Surveys and Monographs, vol.~185,
\newblock American Mathematical Society, Providence, RI, 2013.

\bibitem{KamberTondeur1975}
F.~W.~Kamber and P.~Tondeur,
\newblock {\em Foliated Bundles and Characteristic Classes},
\newblock Lecture Notes in Mathematics, vol.~493,
\newblock Springer, Berlin--New York, 1975.

\bibitem{Karoubi1978}
M.~Karoubi,
\newblock {\em K-Theory: An Introduction},
\newblock Grundlehren der mathematischen Wissenschaften, vol.~226,
\newblock Springer, Berlin--New York, 1978.

\bibitem{Kock2006}
A.~Kock,
\newblock {\em Synthetic Differential Geometry},
\newblock 2nd ed.,
\newblock London Mathematical Society Lecture Note Series, vol.~333,
\newblock Cambridge University Press, Cambridge, 2006.

\bibitem{KrieglMichor1997}
A.~Kriegl and P.~W.~Michor,
\newblock {\em The Convenient Setting of Global Analysis},
\newblock Mathematical Surveys and Monographs, vol.~53,
\newblock American Mathematical Society, Providence, RI, 1997.

\bibitem{MagnotConnections}
J.-P.~Magnot,
\newblock [Exact title of the article on diffeological connections],
\newblock {\em Carpathian J. Math.},
\newblock [volume, year, pages to be completed].

\bibitem{MagnotWatts2017}
J.-P.~Magnot and J.~Watts,
\newblock The diffeology of Milnor's classifying space,
\newblock {\em Topology Appl.} 232 (2017), 189--213.

\bibitem{Magnot-MilnorI}
J.-P. Magnot,
\newblock On spherical Milnor Classifying Spaces I:
differential geometry
\newblock Preprint 2026

\bibitem{Milnor1956I}
J.~Milnor,
\newblock Construction of universal bundles. I,
\newblock {\em Ann. of Math.} (2) 63 (1956), 272--284.

\bibitem{Milnor1956II}
J.~Milnor,
\newblock Construction of universal bundles. II,
\newblock {\em Ann. of Math.} (2) 63 (1956), 430--436.

\bibitem{MoerdijkPronk1997}
I.~Moerdijk and D.~A.~Pronk,
\newblock Orbifolds, sheaves and groupoids,
\newblock {\em K-Theory} 12 (1997), no.~1, 3--21.

\bibitem{Murray1996}
M.~K.~Murray,
\newblock Bundle gerbes,
\newblock {\em J. London Math. Soc.} (2) 54 (1996), no.~2, 403--416.

\bibitem{Neeb2004}
K.-H.~Neeb,
\newblock Abelian extensions of infinite-dimensional Lie groups,
\newblock {\em Travaux Math\'ematiques} 15 (2004), 69--194.

\bibitem{Neeb2007}
K.-H.~Neeb,
\newblock Non-abelian extensions of infinite-dimensional Lie groups,
\newblock {\em Ann. Inst. Fourier (Grenoble)} 57 (2007), no.~1, 209--271.

\bibitem{NikolausSchreiberStevenson2012}
T.~Nikolaus, U.~Schreiber, and D.~Stevenson,
\newblock Principal $\infty$-bundles: general theory,
\newblock {\em J. Homotopy Relat. Struct.} 10 (2015), no.~4, 749--801.

\bibitem{NikolausWaldorf2013}
T.~Nikolaus and K.~Waldorf,
\newblock Four equivalent versions of non-abelian gerbes,
\newblock {\em Pacific J. Math.} 264 (2013), no.~2, 355--420.

\bibitem{Noohi2005}
B.~Noohi,
\newblock Foundations of topological stacks I,
\newblock arXiv:math/0503247.

\bibitem{Noohi2012}
B.~Noohi,
\newblock Homotopy types of topological stacks,
\newblock {\em Adv. Math.} 230 (2012), no.~4--6, 2014--2047.

\bibitem{Pervova2016}
E.~Pervova,
\newblock Diffeological vector pseudo-bundles,
\newblock {\em Topology Appl.} 202 (2016), 269--300.

\bibitem{Pervova2017}
E.~Pervova,
\newblock Diffeological connections on diffeological vector pseudo-bundles,
\newblock arXiv:1611.07694.

\bibitem{Pervova2018}
E.~Pervova,
\newblock Levi-Civita connections on diffeological vector pseudo-bundles,
\newblock arXiv:1701.04988.

\bibitem{PressleySegal1986}
A.~Pressley and G.~Segal,
\newblock {\em Loop Groups},
\newblock Oxford Univ. Press, 1986.

\bibitem{Rector1971}
D.~L.~Rector,
\newblock Steenrod operations in the Eilenberg--Moore spectral sequence,
\newblock {\em Comment. Math. Helv.} 45 (1970), 540--552.


\bibitem{Rosenberg1989}
J.~Rosenberg,
\newblock Continuous-trace algebras from the bundle theoretic point of view,
\newblock {\em J. Austral. Math. Soc. Ser. A} 47 (1989), no.~3, 368--381.

\bibitem{Scalbi2024}
N.~T.~Scalbi,
\newblock Differential cohomology as diffeological homotopy theory,
\newblock arXiv:2408.02593.

\bibitem{Schreiber2013}
U.~Schreiber,
\newblock Differential cohomology in a cohesive $\infty$-topos,
\newblock arXiv:1310.7930.

\bibitem{Segal1968}
G.~Segal,
\newblock Classifying spaces and spectral sequences,
\newblock {\em Inst. Hautes \'Etudes Sci. Publ. Math.} 34 (1968), 105--112.

\bibitem{Souriau1980}
J.-M.~Souriau,
\newblock Groupes diff\'erentiels,
\newblock in {\em Differential Geometrical Methods in Mathematical Physics},
\newblock Lecture Notes in Mathematics, vol.~836,
\newblock Springer, Berlin, 1980, pp.~91--128.

\bibitem{Stacey2011}
A.~Stacey,
\newblock Comparative smootheology,
\newblock {\em Theory Appl. Categ.} 25 (2011), no.~4, 64--117.

\bibitem{Stevenson2000}
D.~Stevenson,
\newblock The geometry of bundle gerbes,
\newblock Ph.D. thesis, University of Adelaide, 2000.

\bibitem{Waldorf2007}
K.~Waldorf,
\newblock More morphisms between bundle gerbes,
\newblock {\em Theory Appl. Categ.} 18 (2007), no.~9, 240--273.

\bibitem{Waldorf2010}
K.~Waldorf,
\newblock A construction of string 2-group models using a transgression-regression technique,
\newblock in {\em Analysis, Geometry and Quantum Field Theory},
\newblock Contemp. Math., vol.~584,
\newblock Amer. Math. Soc., Providence, RI, 2012, pp.~99--115.

\bibitem{Watts2012Thesis}
J.~Watts,
\newblock {\em Diffeologies, Differential Spaces, and Symplectic Geometry},
\newblock Ph.D. thesis, University of Toronto, 2012.

\bibitem{Wockel2008}
C.~Wockel,
\newblock Principal 2-bundles and their classifying spaces,
\newblock arXiv:0803.3692.

\bibitem{Wu2023}
E.~Wu,
\newblock Pushforward and smooth vector pseudo-bundles,
\newblock {\em Pacific J. Math.} 322 (2023), no.~1, 195--224.

\end{thebibliography}
\end{document}
